%% ============================================================
%%  Signal Research Report — BTCUSDT Perpetual Futures
%%  Compile with: pdflatex signal_research_report.tex
%% ============================================================
\documentclass[12pt, a4paper]{article}

\usepackage{amsmath, amssymb, amsthm}
\usepackage{booktabs}
\usepackage{hyperref}
\usepackage{geometry}
\usepackage{xcolor}
\usepackage{array}
\usepackage{multirow}
\usepackage{enumitem}
\usepackage{caption}
\usepackage{float}
\usepackage{parskip}

\geometry{margin=2.5cm}

\hypersetup{
  colorlinks=true,
  linkcolor=blue!60!black,
  urlcolor=blue!60!black
}

\theoremstyle{definition}
\newtheorem{definition}{Definition}[section]
\newtheorem{proposition}{Proposition}[section]

\newcommand{\IC}{\mathrm{IC}}
\newcommand{\E}{\mathbb{E}}
\newcommand{\Var}{\mathrm{Var}}
\newcommand{\Cov}{\mathrm{Cov}}
\newcommand{\sign}{\mathrm{sign}}
\newcommand{\rank}{\mathrm{rank}}
\newcommand{\EMA}{\mathrm{EMA}}

\title{\textbf{Systematic Signal Research on BTCUSDT Perpetual Futures}\\[0.4em]
\large Phase 1: Order Flow \& Arbitrage Signals (A--J)\\[0.2em]
\large Phase 2: Mechanism-Based Screening (D4, N)}
\author{}
\date{May 2026}

\begin{document}
\maketitle
\tableofcontents
\newpage

%% ============================================================
\section{Introduction}
%% ============================================================

This report documents the design, mathematical construction, and rigorous
statistical evaluation of ten predictive signals on the BTCUSDT USDT-margined
perpetual futures contract traded on Binance.

The central question is not merely whether a signal predicts future returns
but whether the gross prediction edge exceeds the round-trip transaction cost.
Every signal is therefore evaluated along two axes:
\begin{enumerate}
  \item \textbf{Statistical significance}: is the measured IC distinguishable
        from zero after accounting for autocorrelation and multiple testing?
  \item \textbf{Economic significance}: does the IC exceed the breakeven
        threshold $\IC^*$ implied by transaction costs?
\end{enumerate}

A key finding that emerges across Phase 1 is that statistically significant
signals may still be economically unviable if they belong to the same
\emph{microstructure family} — short-horizon taker-side order flow mean reversion —
which is competed away by market makers operating with lower costs and faster
infrastructure. Phase 1 therefore broadened from order-flow signals
(Signals C, E\textsuperscript{vpin}, B, F) toward mechanism-driven signals rooted
in \emph{forced} or \emph{arbitrage-compelled} trading behaviour (Signals H1, H2, I).

Phase 2 shifted to five structurally distinct mechanisms: open interest velocity,
liquidation cascades, cross-exchange funding dispersion, and options volatility
regimes. A first-pass IC screen on nine signals identified the DVOL family
(N1--N3) as first-pass survivors; rigorous daily bootstrap validation reduces
this to \textbf{N3 alone} (Deribit 30-day implied volatility $z$-score) as the
dominant survivor --- N1 and N2 do not pass the non-overlapping daily
significance test. D4 (funding$\times\Delta$OI), which appeared promising at
first pass, was killed when evaluated correctly at settlement bars. N3 achieves
$|\IC|/\IC^*_{\text{maker}} = 4.11\times$ at the 24h horizon on non-overlapping
daily returns ($p = 0.039$, block-bootstrap on daily series), the strongest
result of the programme.


%% ============================================================
\section{Data}
%% ============================================================

\subsection{Data Source and Period}

The primary dataset is the BTCUSDT USDT-margined perpetual contract on Binance
Futures (symbol \texttt{BTCUSDT}, FAPI endpoint), sampled at 1-minute
frequency. The full dataset spans January 2022 to May 2026.

Each bar $t$ contains:
\[
  \mathcal{B}_t = \bigl(O_t,\, H_t,\, L_t,\, C_t,\, V_t,\, Q_t,\,
                        V_t^B,\, Q_t^B,\, N_t,\, M_t,\, F_t\bigr)
\]
where:
\begin{itemize}[noitemsep]
  \item $C_t$ — close price (last traded price)
  \item $V_t$ — total base-asset volume (BTC)
  \item $Q_t$ — total quote-asset volume (USDT)
  \item $V_t^B$ — taker buy base volume; taker sell is $V_t^S = V_t - V_t^B$
  \item $Q_t^B$ — taker buy quote volume
  \item $N_t$ — number of trades
  \item $M_t$ — 8-hour funding mark price (constant within each 8h period)
  \item $F_t$ — 8-hour funding rate (applied at 00:00, 08:00, 16:00 UTC)
\end{itemize}

\subsection{Forward Return Targets}

For prediction horizon $h \in \{1\text{m}, 5\text{m}, 15\text{m}, 1\text{h},
4\text{h}, 8\text{h}\}$, the forward log-return is:
\[
  r_t^{(h)} = \log C_{t+h} - \log C_t
\]
All signals are evaluated as predictors of $r_t^{(h)}$.

\subsection{Train / Validation / Test Split}

The dataset is split chronologically (no overlap, no leakage):
\begin{center}
\begin{tabular}{lll}
\toprule
Split & Period & Bars (1m) \\
\midrule
Train & 2023-01-01 to 2023-12-31 & 525,600 \\
Val   & 2024-01-01 to 2024-06-30 & 262,080 \\
Test  & 2024-07-01 to 2024-12-31 & 264,960 \\
\bottomrule
\end{tabular}
\end{center}

\subsection{Transaction Cost Model}

Binance USDT-M perpetual fees for a standard account (VIP0):
\begin{align*}
  c_{\text{taker}} &= 0.040\% \text{ (fee)} + 0.010\% \text{ (slippage)} = 0.050\%
  \text{ per side}\\
  c_{\text{RT}}^{\text{taker}} &= 2 \times 0.050\% = 0.10\% \text{ round-trip}\\
  c_{\text{RT}}^{\text{maker}} &= 2 \times (0.020\% + 0.010\%) = 0.06\%
  \text{ round-trip}
\end{align*}


%% ============================================================
\section{Statistical Framework}
%% ============================================================

\subsection{Spearman Information Coefficient}

\begin{definition}[Spearman IC]
  For signal values $\{s_t\}$ and forward returns $\{r_t^{(h)}\}$, the
  Spearman Information Coefficient is:
  \[
    \IC(s, r^{(h)}) = \rho_S\bigl(\rank(s_t),\, \rank(r_t^{(h)})\bigr)
    = 1 - \frac{6 \sum_{t=1}^{n} d_t^2}{n(n^2-1)}
  \]
  where $d_t = \rank(s_t) - \rank(r_t^{(h)})$. The Spearman IC is preferred
  over Pearson correlation because it is robust to outliers and monotone
  transformations of the signal.
\end{definition}

A negative IC is not intrinsically worse than a positive IC; what matters is
$|\IC| > \IC^*$. Throughout this report, the sign of IC indicates
\emph{direction} (negative means the signal is a contrarian predictor).

\subsection{Permutation Test}

The permutation test assesses $H_0 : \IC = 0$ by constructing the null
distribution empirically:

\begin{enumerate}
  \item Compute observed $\widehat{\IC}$ on the original sample.
  \item For $b = 1,\ldots,B$ (typically $B = 2{,}000$): shuffle $\{s_t\}$
        uniformly at random while keeping $\{r_t^{(h)}\}$ fixed, yielding
        $\widehat{\IC}^{(b)}$.
  \item $p = \frac{1}{B}\sum_{b=1}^B \mathbf{1}\bigl[|\widehat{\IC}^{(b)}|
        \geq |\widehat{\IC}|\bigr]$.
\end{enumerate}

Shuffling the signal while fixing returns preserves the serial structure of
returns (relevant for detecting momentum/mean-reversion in the target) but
destroys any predictive relationship. For large samples ($n > 50{,}000$),
a random subsample of size $50{,}000$ is used for computational efficiency.

\subsection{Circular Block Bootstrap}

Because both signals and returns exhibit autocorrelation, permutation tests
that treat observations as i.i.d.\ are anti-conservative (they overstate
significance). The circular block bootstrap (Politis \& Romano, 1994) provides a
more robust test:

\begin{enumerate}
  \item Choose block length $\ell$ and treat the time series as circular
        (index modulo $n$), so every position is a valid block start.
  \item Draw $\lceil n/\ell \rceil$ block start positions uniformly at random
        with replacement from $\{0, 1, \ldots, n-1\}$; extract contiguous
        blocks of length $\ell$, wrapping around the end if necessary.
  \item Concatenate the resampled blocks and trim to length $n$; compute
        $\widehat{\IC}^{(b)}$ on the resampled series.
  \item Repeat $B$ times.  Report the percentile 95\% confidence interval
        and the $p$-value
        $p = \frac{1}{B}\sum_{b=1}^B \mathbf{1}\bigl[|\widehat{\IC}^{(b)}
        - \widehat{\IC}| \geq |\widehat{\IC}|\bigr]$,
        which shifts the bootstrap distribution to the null
        before comparing (see Section~\ref{sec:h1} for derivation of why
        the unshifted formula gives $p \approx 0.5$ regardless of
        true significance).
\end{enumerate}

The block length $\ell$ is chosen to exceed the signal's autocorrelation
length: $\ell = 30$ bars for 5m signals, $\ell = 60$ bars for 1h signals,
$\ell = 3$ settlements for 8h-settlement signals.

\subsection{Walk-Forward Out-of-Sample Validation}

All signals are evaluated on data outside the training set using a
walk-forward scheme:
\begin{enumerate}
  \item Divide the OOS period (val + test) into windows of $W$ bars
        ($W = 30 \times 1440$ for 1m bars, $W = 90 \times 3$ for settlement bars).
  \item Compute $\IC_k$ within each window $k$ without any re-fitting.
  \item Report mean $\overline{\IC}$, standard deviation $\hat{\sigma}_{\IC}$,
        and the fraction of windows with the same sign as the in-sample IC.
\end{enumerate}

\subsection{Breakeven Information Coefficient}

\begin{proposition}[Breakeven IC]
  Assume a signal $s_t$ is used to trade at horizon $h$, entering when
  $|s_t| > \text{threshold}$ and exiting after $h$ bars. Under a linear
  IC model $r^{(h)}_t = \IC \cdot z_t + \varepsilon_t$ where $z_t$ is the
  standardised signal and $\varepsilon_t \sim \mathcal{N}(0, \sigma_r^2)$,
  the expected gross profit per trade is $\IC \cdot \sigma_r \cdot \E[|z_t|]$.
  For a half-normal signal, $\E[|z_t|] = \sqrt{2/\pi}$. The breakeven IC is:
  \[
    \IC^* = \frac{c_{\text{RT}}}{\sigma_r^{(h)} \cdot \sqrt{2/\pi}}
  \]
  A signal is economically viable only if $|\IC| > \IC^*$.

  \textbf{Caveat — conditional entry}: the formula uses the \emph{unconditional}
  expectation $\E[|z_t|] = \sqrt{2/\pi}$.  If the strategy only trades when the
  signal is extreme (e.g., $|z_t| > c$ for some threshold $c > 0$), the correct
  denominator becomes the conditional expectation
  $\E[|z_t| \mid |z_t| > c] > \sqrt{2/\pi}$, which \emph{lowers} the breakeven
  IC (the signal generates more profit per trade in the tails).  Conversely,
  if trades are executed on every bar regardless of signal strength, the
  unconditional formula applies.  All breakeven IC values reported in this
  document assume the unconditional case; they are therefore conservative
  estimates of tradability for threshold-entry strategies.
\end{proposition}

\begin{table}[H]
\centering
\caption{Breakeven IC thresholds by horizon}
\label{tab:breakeven}
\begin{tabular}{lrrrrr}
\toprule
Horizon & $\sigma_r$ & $\IC^*_{\text{taker}}$ & $\IC^*_{\text{maker}}$ \\
\midrule
1m  & 0.062\% & 2.032 & 1.219 \\
5m  & 0.135\% & 0.926 & 0.556 \\
15m & 0.229\% & 0.548 & 0.329 \\
1h  & 0.449\% & 0.279 & 0.168 \\
4h  & 0.902\% & 0.139 & 0.083 \\
8h  & 1.284\% & 0.098 & 0.059 \\
\bottomrule
\end{tabular}
\end{table}

The key insight from Table~\ref{tab:breakeven}: at short horizons the
breakeven IC is extremely high because $\sigma_r$ is small relative to the
fixed round-trip cost. Profitable short-horizon trading requires either
near-perfect prediction or maker-side execution.


%% ============================================================
\section{Signal C: CVD Divergence}
%% ============================================================

\subsection{Economic Mechanism}

Cumulative Volume Delta (CVD) measures the net directional pressure from
aggressive (taker) order flow. The \emph{divergence} between this flow
and realised price returns captures a crowding imbalance: if aggressive
buyers have been active but price has not risen correspondingly, it implies
that passive sellers (market makers) have been absorbing the flow at the
current price. Once the aggressive buying exhausts, price reverts.

This is an \emph{order flow absorption} hypothesis: the signal fires when
liquidity providers have absorbed directional flow without yielding ground.

\subsection{Signal Construction}

\begin{definition}[CVD Divergence]
  For window $w = 60$ bars and rolling history $H = 300$ bars:
  \begin{align}
    \delta_t^{\text{CVD}} &= \frac{\sum_{s=t-w+1}^{t}(V_s^B - V_s^S)}
                              {\sum_{s=t-w+1}^{t} V_s}
                           - \log\frac{C_t}{C_{t-w}}
  \end{align}
  The first term is the net buy fraction over the past hour; the second is
  the log price return over the same window. A large positive value means
  heavy buying without commensurate price appreciation.

  The trading signal is the \textbf{percentile rank} of $\delta_t^{\text{CVD}}$
  within a rolling window of 300 recent values:
  \[
    \psi_t^C = \frac{1}{|H_t|}\sum_{h \in H_t}
               \mathbf{1}[\delta_h^{\text{CVD}} < \delta_t^{\text{CVD}}]
  \]
  Trade rules: SHORT when $\psi_t^C \geq 0.80$; LONG when $\psi_t^C \leq 0.20$.
\end{definition}

The construction is $O(1)$ per bar by maintaining running sums for the buy,
sell, and volume buffers of length $w+1$, and a sorted mirror list for the
rolling percentile rank (inserting/evicting in $O(\log H)$ via binary search).

\subsection{Results}

\begin{table}[H]
\centering
\caption{Signal C: CVD Divergence — hypothesis test results}
\begin{tabular}{lrrrr}
\toprule
Test & Horizon & IC & $p$-value & $|\IC|/\IC^*$ \\
\midrule
Permutation & 5m  & $-0.040$ & $<0.001$ & $0.043\times$ \\
Permutation & 15m & $-0.049$ & $<0.001$ & $0.089\times$ \\
Permutation & 1h  & $-0.056$ & $<0.001$ & $0.200\times$ \\
Block bootstrap ($\ell=60$) & 1h & $-0.056$ & $<0.001$ & $0.200\times$ \\
\midrule
Walk-forward OOS (336 windows) & 1h & $-0.039$ & -- & -- \\
\quad Frac.\ same direction & & & & $94.9\%$ \\
\bottomrule
\end{tabular}
\end{table}

\textbf{Statistical verdict}: highly significant at all horizons ($p < 0.001$)
under both permutation and block bootstrap; the OOS IC is directionally stable
($94.9\%$ of 30-day windows). Granger causality confirms that lagged CVD adds
predictive power beyond an autoregressive model of returns at lags 2--5
($p < 0.005$).

\textbf{Economic verdict}: the best ratio $|\IC|/\IC^* = 0.20\times$ at the
1h horizon. The gross edge ($\approx 0.02\%$ per trade) is five times smaller
than the round-trip taker cost ($0.10\%$). The signal belongs to the
well-documented family of \emph{microstructure mean reversion} already
exploited by market makers, leaving no edge for a taker strategy.


%% ============================================================
\section{Signal E: VPIN Regime Filter}
%% ============================================================

\subsection{Economic Mechanism}

VPIN (Volume-synchronised Probability of Informed Trading) approximates
the probability that the current order flow is \emph{informed}. In a
high-VPIN regime, one side of the order book is being systematically
picked off by informed traders, and price moves are driven by information
rather than noise. Critically, in such regimes, a market maker who quotes
passively bears high adverse selection risk: the other party \emph{knows}
which way the price will move.

The hypothesis tested here is the reverse: \textbf{low VPIN} (symmetric,
uninformed flow) is a prerequisite for the CVD mean-reversion effect. When
VPIN is high, the directional flow is genuinely informative and
mean-reversion fails.

\subsection{Signal Construction}

\begin{definition}[VPIN Approximation]
  The per-bar flow imbalance is:
  \[
    \iota_t = \frac{|V_t^B - V_t^S|}{V_t}
  \]
  The VPIN over a rolling window of $W = 50$ bars is:
  \[
    \text{VPIN}_t = \frac{1}{W}\sum_{s=t-W+1}^{t} \iota_s
  \]
  Maintained $O(1)$ with a running sum and a deque of length $W$.
\end{definition}

The strategy (Signal E in the codebase as \texttt{VPINRegimeStrategy}) enters
only when VPIN is in the bottom $40\%$ of its 300-bar rolling distribution
(chaotic, symmetric flow regime) and uses CVD divergence for direction.
Position size scales with $(1 - \psi_t^{\text{VPIN}})$, the complement of
the VPIN percentile rank, so larger sizes are allocated in more chaotic regimes.

\subsection{Results}

\begin{table}[H]
\centering
\caption{Signal E: VPIN as volatility regime detector}
\begin{tabular}{llr}
\toprule
Test & Target & Result \\
\midrule
Block bootstrap ($\ell=50$) & $|\IC|(\text{VPIN}, |r^{(1h)}|) = 0.355$ & $p < 0.001$ \\
Kolmogorov--Smirnov & Low-VPIN vs high-VPIN $|r^{(1h)}|$ & $p < 10^{-300}$ \\
Mann--Whitney U & Low-VPIN $|r^{(1h)}|$ vs high-VPIN & $p < 10^{-300}$ \\
\midrule
Low VPIN: mean $|r^{(1h)}|$ & & $0.351\%$ \\
High VPIN: mean $|r^{(1h)}|$ & & $0.175\%$ \\
VPIN as vol classifier (F1) & & $0.465$ \\
OFI IC in low-VPIN regime & & $-0.034$ \\
OFI IC in high-VPIN regime & & $-0.024$ \\
\bottomrule
\end{tabular}
\end{table}

\textbf{Key finding:} VPIN is an exceptionally strong predictor of
\emph{volatility} ($|\IC| = 0.355$, 51$\sigma$ from null). Low-VPIN bars have
twice the realised volatility of high-VPIN bars. However, VPIN is not a useful
\emph{directional} predictor; the regime filter improves the CVD signal but
does not make it economically viable. The best backtest Sharpe ratio across all
parameter configurations is $-8.8$.


%% ============================================================
\section{Signal A: Funding Rate Curvature}
%% ============================================================

\subsection{Economic Mechanism}

The funding rate $F_t$ is the 8-hourly transfer between long and short
position holders, proportional to the deviation of the perpetual price from
the spot index. A sustained positive funding implies crowded longs. The
\emph{curvature} $d^2F/dt^2$ measures whether the crowding is
\textbf{accelerating} (more extreme) or \textbf{decelerating} (peak crowding,
about to reverse).

The exhaustion hypothesis: when funding is large in magnitude but is
\emph{decelerating} (second derivative has opposite sign), sentiment is
reaching its maximum and a reversal is likely.

\subsection{Signal Construction}

Let $\{F_{k}\}_{k=0}^{K}$ be the sequence of 8-hourly funding rates at
settlement timestamps. Define:
\begin{align}
  d^{(1)}_k &= F_k - F_{k-1} \quad \text{(first difference)}\\
  d^{(2)}_k &= d^{(1)}_k - d^{(1)}_{k-1} = F_k - 2F_{k-1} + F_{k-2}
               \quad \text{(second difference, raw curvature)}\\
  E_k &= -F_k \cdot d^{(2)}_k \quad \text{(exhaustion signal)}
\end{align}
$E_k > 0$ when: (i) positive funding decelerating, or (ii) negative funding
decelerating. Both represent sentiment reversal.

The signal is forward-filled to the 1-minute grid:
$e_t = E_k$ for $t \in [\tau_k, \tau_{k+1})$.

\subsection{Results}

\begin{table}[H]
\centering
\caption{Signal A: Funding Rate Curvature}
\begin{tabular}{lrrr}
\toprule
Test & Horizon & IC & $p$-value \\
\midrule
Permutation & 8h (settle) & $-0.033$ & $0.268$ \\
Block bootstrap ($\ell=3$) & 8h (settle) & $-0.033$ & $0.254$ \\
Walk-forward OOS (10 windows) & 8h & $+0.018$ & -- \\
\midrule
Conditional (top 50\% $|F|$) & 8h & $+0.009$ & -- \\
Conditional (top 33\% $|F|$) & 8h & $+0.033$ & -- \\
\bottomrule
\end{tabular}
\end{table}

\textbf{Verdict}: The exhaustion signal is statistically indistinguishable
from noise ($p = 0.268$ permutation, $p = 0.254$ block bootstrap). The OOS
IC reverses sign ($+0.018$ vs $-0.033$ in-sample), a classic overfitting
signature. Only 1,095 settlement events exist in the training window,
providing insufficient power to detect a weak effect. \textbf{Rejected.}


%% ============================================================
\section{Signal B: Informed Flow}
%% ============================================================

\subsection{Economic Mechanism}

Informed traders who possess private information typically execute in larger
trade sizes than uninformed retail flow. When the average trade size in the
current bar is unusually large \emph{and} directionally aligned with the
quote-OFI (order flow imbalance in dollar terms), the combined signal
proxies for institutional/informed order placement.

However, the empirical finding is a \emph{reversal} signal, not a
continuation signal: large informed-looking trades predict \emph{negative}
future returns. This is consistent with the interpretation that
``informed traders'' in this context are primarily market makers and stat-arb
desks providing liquidity against retail momentum, not trend-following funds.

\subsection{Signal Construction}

\begin{definition}[Informed Flow Signal]
  \begin{align}
    \bar{q}_t &= \frac{Q_t}{N_t} \quad \text{(average trade size in USDT)}\\
    \text{OFI}^Q_t &= \frac{Q_t^B}{Q_t} - 0.5 \quad \text{(quote-OFI, centred)}\\
    z_t^{(60)} &= \frac{\bar{q}_t - \mu_t^{(60)}}{\sigma_t^{(60)}}
    \quad \text{(60-bar rolling z-score of trade size)}\\
    \text{IF}_t^{(60)} &= z_t^{(60)} \cdot \text{OFI}^Q_t
    \quad \text{(informed flow)}
  \end{align}
  The smoothed version $\text{IF}_t^{(5)} = \text{MA}_5(\text{IF}_t^{(60)})$
  is also tested.
\end{definition}

\subsection{Results}

\begin{table}[H]
\centering
\caption{Signal B: Informed Flow — hypothesis tests}
\begin{tabular}{lrrrr}
\toprule
Test & Horizon & IC & $p$-value & $\sigma$ \\
\midrule
Permutation & 5m  & $-0.020$ & $<0.001$ & $-4.3$ \\
Permutation & 15m & $-0.014$ & $0.003$  & $-3.1$ \\
Permutation & 1h  & $-0.010$ & $0.023$  & $-2.3$ \\
Permutation (smoothed $s=5$) & 5m & $-0.026$ & $<0.001$ & $-6.0$ \\
Block bootstrap ($\ell=30$) & 5m & $-0.020$ & $<0.001$ & $-4.4$ \\
Block bootstrap ($\ell=60$) & 1h & $-0.010$ & $0.023$  & $-2.3$ \\
Walk-forward OOS (336 windows) & 5m & $-0.013$ & -- & -- \\
\midrule
\multicolumn{5}{l}{\textit{Kurtosis of trade size vs volatility:}}\\
$\IC(\kappa_t^{300}, |r_t^{(1h)}|)$ & 1h & $-0.101$ & $<0.001$ & $-13.3$ \\
\bottomrule
\end{tabular}
\end{table}

\textbf{Additional finding — trade-size kurtosis as volatility predictor:}
The 300-bar rolling excess kurtosis of $\bar{q}_t$:
\[
  \kappa_t^{(300)} = \frac{\E[(\bar{q} - \mu)^4]}{\sigma^4} - 3
\]
predicts $|r_t^{(1h)}|$ with $\IC = -0.101$ ($p < 0.001$, $-13.3\sigma$).
\emph{High kurtosis} (heavy-tailed trade size distribution, i.e., clustering
of very large trades) is associated with \emph{lower} future volatility.
This suggests that large informed traders front-run volatility events,
causing price discovery to be more efficient (less residual volatility)
after bursty flow.

\textbf{Economic verdict:} The directional $\IC \approx 0.020$ at 5m is
$0.022\times$ the taker breakeven of $0.926$. Despite strong statistical
significance, the signal is sub-breakeven by more than an order of magnitude.
The kurtosis-volatility finding is more interesting as a regime filter or
for volatility targeting, not for directional alpha.


%% ============================================================
\section{Signal F: Directional VPIN Asymmetry}
%% ============================================================

\subsection{Economic Mechanism}

Standard VPIN is symmetric: it measures $|V^B - V^S|/V$ and discards
the \emph{direction} of the imbalance. The directional VPIN asymmetry
instead maintains separate EMA-smoothed buy and sell intensities. When
the buy side has been consistently dominant over 50 bars, the available
ask-side liquidity has been absorbed; market makers widen their spreads
and price reverts.

\subsection{Signal Construction}

\begin{definition}[Directional VPIN Asymmetry]
  \begin{align}
    b_t &= \frac{V_t^B}{V_t} \quad \in [0,1] \quad \text{(buy fraction)}\\
    \text{VPIN}^B_t &= \EMA_{50}(b_t)\\
    \text{VPIN}^S_t &= \EMA_{50}(1 - b_t) = 1 - \text{VPIN}^B_t\\
    A_t &= \text{VPIN}^B_t - \text{VPIN}^S_t = 2\,\EMA_{50}(b_t) - 1
  \end{align}
  Since $\text{VPIN}^S_t = 1 - \text{VPIN}^B_t$, the asymmetry $A_t$ is
  a bijective transformation of $\EMA_{50}(b_t)$, which equals
  $\EMA_{50}(\text{OFI}_t)$ up to a constant shift. The z-score normalised
  version:
  \[
    \tilde{A}_t = \frac{A_t - \mu_t^{(300)}}{\sigma_t^{(300)}}
  \]
\end{definition}

\textbf{Important structural observation}: because $b_t + (1-b_t) = 1$
identically, the asymmetry $A_t$ is mathematically equivalent to a centred
and scaled version of $\EMA_{50}(\text{OFI}_t)$. Signal F reduces to the
same underlying quantity as Signal B's smoothed OFI, confirming that the
entire family of order flow signals (CVD, OFI, informed flow, directional
VPIN) measures the same microstructure phenomenon.

\subsection{Results}

\begin{table}[H]
\centering
\caption{Signal F: Directional VPIN Asymmetry — hypothesis tests}
\begin{tabular}{lrrrr}
\toprule
Test & Horizon & IC & $p$-value & $\sigma$ \\
\midrule
Permutation & 5m  & $-0.040$ & $<0.001$ & $-8.8$ \\
Permutation & 15m & $-0.050$ & $<0.001$ & $-11.0$ \\
Permutation & 1h  & $-0.037$ & $<0.001$ & $-7.9$ \\
Block bootstrap CI (95\%) & 5m & $[-0.046, -0.036]$ & -- & -- \\
Walk-forward OOS (48 windows) & 5m & $-0.020$ & -- & -- \\
\quad Frac.\ same direction & & & & $91.7\%$ \\
\bottomrule
\end{tabular}
\end{table}

\textbf{Lead-lag profile:}
\[
  \IC(A_t,\, r_{t+\ell}^{(1m)}) \approx -0.026 + 0.030 \cdot \mathbf{1}[\ell > 0] \cdot \exp(-\ell/20)
\]
The signal decays from $-0.026$ at lag 0 to approximately $-0.002$ at lag 60,
with a half-life of roughly 20 bars.

\textbf{Volatility conditioning:}
\[
  \IC(A_t, r^{(5m)} \mid \text{high vol}) = -0.050, \quad
  \IC(A_t, r^{(5m)} \mid \text{low vol})  = -0.029
\]
The signal is $70\%$ stronger in high-volatility regimes.

\textbf{Economic verdict:} Best ratio $0.131\times$ taker breakeven at 1h.
Statistically robust and directionally stable OOS, but economically
sub-breakeven by $7.6\times$ at taker costs.


%% ============================================================
\section{Signal H1: Basis Dislocation}
\label{sec:h1}
%% ============================================================

\subsection{Economic Mechanism}

In a USDT-margined perpetual, the funding rate is computed every 8 hours
based on the \emph{premium index} — the deviation of the perpetual mark price
from the spot index. When the perpetual trades at a premium to its funding
mark (set 8 hours ago), two deterministic forces converge it:
\begin{enumerate}
  \item \textbf{Funding arbitrage}: the positive funding payment causes
        longs to pay shorts, reducing the incentive to hold long positions
        and increasing the incentive to short.
  \item \textbf{Post-settlement re-anchoring}: at each 8h settlement, the
        mark resets to the current price, wiping the accumulated basis.
        The \emph{direction} of return in the hour after settlement tends
        to reverse the accumulated basis.
\end{enumerate}

This is a \emph{scheduled, deterministic arbitrage pressure} distinct from
the statistical patterns driving Signals C, B, and F.

\subsection{Signal Construction}

\begin{definition}[Basis Dislocation]
  \[
    b_t = \frac{C_t - M_t}{M_t}
  \]
  where $M_t$ is the 8-hour funding mark price (constant within each 8h period,
  updated at funding settlement). The normalised version:
  \[
    \tilde{b}_t = \frac{b_t - \mu_t^{(480)}}{\sigma_t^{(480)}}
  \]
  where the rolling statistics use only the 480 bars ($= 8$ hours) preceding $t$.
\end{definition}

The signal is evaluated only at \textbf{settlement bars} (timestamps
$\tau_k$ where $F_{\tau_k}$ is applied) to avoid look-ahead contamination
from the next funding mark.

\subsection{Caution: Lookahead Bias in Naive Implementations}

A critical pitfall encountered during this research: using
$r_{t-1}^{(h)}$ (the forward $h$-horizon return starting from $t-1$) as a
``past direction'' proxy is catastrophically biased. For $h = 8\text{h}$
(480 bars) and a 1-bar shift:
\[
  \text{Overlap} = \frac{480 - 1}{480} = 99.8\%
\]
The ``past'' return almost entirely overlaps with the prediction target,
creating spurious ICs of $|\IC| \approx 0.71$ that vanish ($|\IC| < 0.03$)
once corrected by using $C_t / C_{t-480} - 1$ as the true backward return.

\subsection{Results}

\begin{table}[H]
\centering
\caption{Signal H1: statistical tests (training set, 2023, $n=1{,}095$ settlement bars)}
\begin{tabular}{lrrr}
\toprule
Test & Horizon & IC & $p$-value \\
\midrule
Permutation ($N=2000$) & 1h & $-0.170$ & $0.021$ \\
Permutation ($N=2000$) & 8h & $+0.019$ & $0.792$ \\
Block bootstrap ($\ell=3$, $B=1000$) 95\% CI & 1h & $[-0.326, -0.027]$ & $0.028$ \\
\midrule
Walk-forward OOS ($W=270$ settlements, 10 windows) & 1h & $-0.029$ & -- \\
\quad Fraction same direction & & & $70.0\%$ \\
\bottomrule
\end{tabular}
\end{table}

\begin{table}[H]
\centering
\caption{Signal H1: OOS performance by calendar period}
\label{tab:h1_oos}
\begin{tabular}{lrrr}
\toprule
Period & $n$ & IC (1h) & Ratio (maker) \\
\midrule
Train 2023         & 1,095 & $-0.170$ & $1.02\times$ \\
Val   2024-H1      & 546   & $+0.008$ & $0.06\times$ \\
Test  2024-H2      & 552   & $-0.044$ & $0.33\times$ \\
\bottomrule
\end{tabular}
\end{table}

\begin{table}[H]
\centering
\caption{Signal H1: breakeven analysis (1h horizon, training)}
\begin{tabular}{lrrrr}
\toprule
Horizon & $\IC^*_{\text{taker}}$ & $\IC^*_{\text{maker}}$ & $|\IC|$ & Ratio (maker) \\
\midrule
1h & $0.277$ & $0.166$ & $0.170$ & $\mathbf{1.02\times}$ \\
8h & $0.098$ & $0.059$ & $0.062$ & $1.05\times$ \\
\bottomrule
\end{tabular}
\end{table}

\textbf{Statistical verdict (training)}: the basis signal is significant at the 5\%
level under both the permutation test ($p = 0.021$) and the block bootstrap
($p = 0.028$, 95\% CI $[-0.326, -0.027]$ excludes zero). These two tests are
now consistent: an earlier draft reported $p = 0.464$ from the bootstrap, which
was an error — the test was checking whether $|\hat{\rho}_{\text{boot}}| \geq |\hat{\rho}|$,
but because the bootstrap distribution is centred near $\hat{\rho}$ rather than
near zero, $\approx$50\% of samples always pass regardless of the CI. The correct
test shifts the bootstrap distribution to the null before computing the tail
probability: $p = \Pr\bigl(|\hat{\rho}^* - \hat{\rho}| \geq |\hat{\rho}|\bigr)$.

\begin{table}[H]
\centering
\small
\caption{Signal H1 regime stability via premium index proxy (2022--2026)$^\dagger$}
\label{tab:h1_regime}
\begin{tabular}{lrrrrr}
\toprule
Period & $n$ & IC (1h) & Ratio 1h & IC (8h) & Ratio 8h \\
\midrule
2022             & 1,095 & $-0.023$ & $0.21\times$ & $-0.085$ & $2.09\times$ \\
2023 (train)     & 1,095 & $-0.054$ & $0.33\times$ & $-0.021$ & $0.36\times$ \\
\midrule
2024-H1 (OOS)    &   546 & $+0.071$ & $0.50\times^\ddagger$ & $+0.005$ & $0.09\times$ \\
2024-H2 (OOS)    &   552 & $-0.010$ & $0.08\times$ & $+0.001$ & $0.02\times$ \\
2025-H1 (OOS)    &   543 & $-0.067$ & $0.47\times$ & $-0.053$ & $\mathbf{1.03\times}$ \\
2025-H2 (OOS)    &   552 & $-0.024$ & $0.14\times$ & $-0.001$ & $0.02\times$ \\
2026 YTD (OOS)   &   395 & $-0.046$ & $0.30\times$ & $-0.008$ & $0.14\times$ \\
\midrule
OOS agg.\ 2024+  & 2,588 & $-0.024$ & -- & -- & -- \\
Train 2022--23   & 2,190 & $-0.025$ & -- & -- & -- \\
\bottomrule
\end{tabular}
\par\smallskip
\raggedright
\footnotesize
$^\dagger$~Signal proxy: $\mathtt{basis\_pct} = (\text{mark} - \text{index})/\text{index}$
from premium index data --- same mechanism as H1 but mark vs.\ real-time spot index
rather than close vs.\ 8h funding mark; ratios relative to maker breakeven.\\
$^\ddagger$~IC sign inverted during Jan--Jun 2024 BTC bull rally (basis expanded persistently
rather than mean-reverting); see OOS verdict below.
\end{table}

\textbf{OOS verdict}: regime stability is confirmed over 2022--2026
(Table~\ref{tab:h1_regime}).  The 2024-H1 IC inversion ($+0.071$ via premium
index proxy) was a \emph{bull-market anomaly}: during the Jan--Jun 2024 rally,
basis persistently expanded rather than mean-reverting, so the funding
arbitrage failed to converge.  Outside this window, the signal maintains
the expected negative sign in every half-year block, including 2025 and 2026.
In 2025-H1 the 8h IC recovers to $-0.053$ ($\mathbf{1.03\times}$ maker
breakeven) --- the first OOS window to clear the cost threshold.  The
aggregate OOS IC over 2024--2026 is $-0.024$ versus train IC $-0.025$:
near-identical, confirming the mechanism is not regime-specific.
Walk-forward OOS (26 windows, 90 days each, 30-day step) shows the correct
sign in 62\% of windows, consistent with a small but persistent edge.

\textbf{Regime conditioning (in-sample illusion)}: conditioning on
$|b_t| > 75\text{th}$ percentile inflates the training IC to $-0.347$
($2.44\times$ maker).  However, this conditioning produces
$\mathrm{IC} \approx 0$ in both 2024 periods.  The extreme-basis
amplification is \emph{spurious} — it captures a 2023-specific regime
and does not generalise.

\textbf{Return timing}: the 15 minutes \emph{before} settlement show
$\mathrm{IC}(b_t, r_{t-15:t-1}) = -0.147$ (training), comparable to the
post-settlement 1h IC.  This implies informed participants are already
repositioning before the funding payment crystallises.  One cannot exploit
this observation directly (basis $b_t$ is only known at $t$), but it suggests
that the premium index mark-vs-index data (Signal J) could serve as a
forward-looking proxy.

\textbf{H1 + H2 combination (OOS result)}:

\begin{table}[H]
\centering
\small
\caption{H1 + H2 equal-weight $z$-score combination: OOS performance (1h horizon)}
\label{tab:h1h2_oos}
\begin{tabular}{lrrrrr}
\toprule
Period & $n$ & IC (H1) & IC (H2) & IC (Combo) & Ratio (maker) \\
\midrule
Train 2023$^\dagger$ &   185 & $-0.170$ & $-0.178$ & $-0.225$ & $1.30\times$ \\
Val   2024-H1        &   546 & $+0.008$ & $+0.047$ & $+0.021$ & $0.15\times$ \\
Test  2024-H2        &   552 & $-0.044$ & $-0.096$ & $-0.096$ & $0.71\times$ \\
\midrule
OOS   2024           & 1,098 & $-0.018$ & $-0.024$ & $-0.038$ & $\mathbf{0.28\times}$ \\
\bottomrule
\end{tabular}
\par\smallskip\raggedright\footnotesize
$^\dagger$~\texttt{basis} is non-NaN for only 185 of 1095 settlement bars in
the 2023 training window of ofi.parquet; the training IC figure is based on this
sparse sample and should be treated with caution.
\end{table}

The combination does \emph{not} rescue OOS performance.  In 2024-H1 both
signals invert to the wrong sign simultaneously --- H1 ($+0.008$) and
H2 ($+0.047$) --- confirming they share the same regime failure.  Low
correlation $\rho = 0.12$ (stable at $0.09$ in both OOS halves) does not
prevent co-directional collapse when the underlying mechanism breaks down.
Walk-forward OOS gives the combo 80\% correct direction versus 70\% for
H1 alone, a marginal stability gain, but with mean IC $= -0.041$ across
10 windows --- far below the maker breakeven of ${\approx}0.15$.
The optimal OOS weight is $0.8\,\text{H1} + 0.2\,\text{H2}$ ($0.29\times$
maker).  \textbf{The combination does not change the execution case}: the
OOS edge is real, below-breakeven, and H2 contributes modest stability
rather than independent alpha.

\textbf{Execution verdict}: \emph{mechanism is real; current execution
structure is not viable.}  Pre-settlement entry (the natural maker approach)
is fatally compromised by two compounding problems.  First, the IC
\emph{inverts} the moment entry is moved before settlement: entering at
$t-1$ produces IC $= +0.097$ (wrong sign), and the inversion persists at all
pre-settlement offsets tested ($t-2$ through $t-30$).  The market front-runs
the settlement: price rises with positive basis in the minutes before the
funding payment rather than after it.  Second, adverse-selection rates are
severe: 91\% of maker fills at $t-1$ arrive on bars that move against the
signal, falling to 56\% at $t-15$ (random would be 50\%).  At-settlement
entry ($k=0$) recovers the correct IC direction ($-0.018$, $0.13\times$
maker) but requires a taker order, which is sub-breakeven by $7\times$.
The sole interesting subset is the 00:00 UTC settlement (IC $= -0.090$,
$0.67\times$ maker OOS) versus near-zero IC at 08:00 and 16:00 --- but this
is still below the maker cost threshold.  \textbf{H1 is not viable as a
standalone single-exchange strategy and will not be developed further.}
The structural insight --- that basis mean-reversion is concentrated at a
specific exchange and settlement hour --- motivates the cross-exchange
relative-basis direction (see Next Steps).

\textbf{Note on the basis variable}: in the processed dataset, $M_t$ is the
8-hour funding mark (a 30-minute TWAP of the perpetual, computed at settlement
and held constant). It differs from the real-time 30-min rolling TWAP mark
used for continuous margin calculations. A cleaner version of this signal,
using the real-time mark vs.\ the cross-exchange spot index, is the subject
of Signal J (premium index download, Section~\ref{sec:premium}).


%% ============================================================
\section{Signal H2: Funding Carry as Directional Signal}
%% ============================================================

\subsection{Economic Mechanism}

A positive funding rate implies that longs are paying shorts every 8 hours.
This creates a carry cost for long holders proportional to their position
size. As funding persists, marginal longs exit (the trade is less attractive),
and shorts receive an ongoing subsidy that incentivises more short positions.
The prediction is that \emph{high positive funding} should be associated with
downward price pressure over the next 8h as the long-side unwinds.

\subsection{Signal Construction}

\[
  \phi_t = F_t, \qquad
  \hat\phi_t = \frac{1}{|H|}\sum_{k \in H} \mathbf{1}[F_k < F_t]
  \quad \text{(percentile rank over 300 settlements)}
\]

Trade rule: SHORT when $\hat\phi_t \geq 0.90$; LONG when $\hat\phi_t \leq 0.10$.

\subsection{Results}

\begin{table}[H]
\centering
\caption{Signal H2: Funding Carry}
\begin{tabular}{lrrr}
\toprule
Test & Horizon & IC & $p$-value \\
\midrule
Permutation ($N=2000$) & 1h (settle) & $-0.067$ & -- \\
Permutation ($N=2000$) & 8h (settle) & $-0.022$ & $0.471$ \\
Walk-forward OOS (10 windows) & 8h & $-0.038$ & -- \\
\quad Fraction correct direction & & & $80.0\%$ \\
\midrule
Extreme only (top/bot 20\%) & 8h & $-0.046$ & -- \\
\bottomrule
\end{tabular}
\end{table}

\textbf{Verdict}: The 8h funding carry signal is not statistically significant
($p = 0.471$). The 10-window OOS evidence shows 80\% directional consistency,
but with only 10 windows this is not conclusive. The ratio
$|\IC|/\IC^*_{\text{taker}} = 0.023/0.098 = 0.23\times$. \textbf{Sub-breakeven.}

The 1h IC at settlement bars ($-0.067$) is the most interesting result and
warrants future investigation with higher-frequency settlement data (e.g.,
exchanges with hourly funding settlements).


%% ============================================================
\section{Signal I: Crowded-Trade Exhaustion}
%% ============================================================

\subsection{Economic Mechanism}

A position holder is \emph{trapped} when:
(i) they are paying carry (positive funding for a long), and
(ii) the position is losing money (price has fallen since they entered).
Such holders face mounting pressure to exit: ongoing funding payments reduce
the carrying capacity, and the unrealised loss compounds the psychological
and risk-management pressure. The eventual exit creates a price cascade.

\subsection{Signal Construction}

\begin{definition}[Trapped Positioning Score]
  \[
    P_t^{\text{past}} = \frac{C_t - C_{t-480}}{C_{t-480}}
    \quad \text{(true past 8h price return)}
  \]
  \[
    \tau_t = \sign(F_t) \cdot \bigl(-\sign(P_t^{\text{past}})\bigr)
  \]
  $\tau_t = +1$ when longs are trapped (positive funding, falling price) or
  shorts are trapped (negative funding, rising price).
  $\tau_t = -1$ when positions are profitable and carry is favourable.
\end{definition}

\subsection{Lookahead Bias Analysis}

The natural implementation mistake is:
\[
  \text{WRONG: } \tau_t = \sign(F_t) \cdot (-\sign(r_{t-1}^{(8h)}))
\]
Since $r_t^{(8h)} = \log C_{t+480} - \log C_t$ is a \emph{forward} return,
$r_{t-1}^{(8h)}$ covers the interval $[t-1, t+479]$, which has
$479/480 = 99.8\%$ temporal overlap with $r_t^{(8h)}$. This creates a
spurious $\IC = -0.71$ that evaporates completely after correction.

The correct implementation uses $C_t/C_{t-480} - 1$:
\begin{center}
\begin{tabular}{lrr}
\toprule
Implementation & IC($\tau_t$, $r_t^{(8h)}$) & $p$-value \\
\midrule
Naive (shift 1 bar) & $-0.710$ & $< 0.001$ \\
Corrected (true past 8h) & $+0.027$ & $0.348$ \\
\bottomrule
\end{tabular}
\end{center}

\textbf{Verdict}: After correcting the lookahead bias, the trapped positioning
signal is statistically indistinguishable from zero. The hypothesis is
interesting economically but not supported in the data over this sample.

\subsection{Joint Extreme: Funding and Basis Jointly Aligned}

A combined signal uses both the funding level and the current basis:
\[
  J_t = \tilde{F}_t \cdot \tilde{b}_t
\]
where $\tilde{F}_t$ is the standardised funding rate and $\tilde{b}_t$ is
the standardised basis. When both are positive (longs crowded and price
above mark), or both negative (shorts crowded and price below mark), the
crowding signal is reinforced.

Results at 1m level: $\IC = -0.040$, $p = 0.001$, $-3.6\sigma$ (permutation),
but $|\IC|/\IC^*_{\text{taker}}(1h) = 0.014\times$. The statistical signal
is real but economically negligible.


%% ============================================================
\section{Signal J: Continuous Mark-Index Basis (Premium Index)}
\label{sec:premium}
%% ============================================================

\subsection{Mechanism}

The premium index basis $b^{\text{prem}}_t = (M_t - I_t)/I_t$ is the
real-time input to the Binance funding rate TWAP, where $M_t$ is the 1-minute
mark price and $I_t$ is the cross-exchange spot composite index.  When
$b^{\text{prem}}_t > 0$, long holders implicitly accumulate a carry cost not
yet crystallised; index arbitrageurs can simultaneously short the perpetual
and buy spot to earn a risk-free spread, exerting downward pressure on $M_t$.
The hypothesis is that an elevated (or depressed) basis predicts a partial
return reversal as this arbitrage pressure resolves continuously.

This is the \emph{continuous-time} analogue of Signal H1 (Section~\ref{sec:h1}).
The key difference is that H1 operates only at the three daily settlement bars
(00:00, 08:00, 16:00 UTC) where funding is mechanically certain; Signal J
operates on all 1-minute bars ($n \approx 525{,}000$ training observations).

\subsection{Signal Construction}

Data source: \texttt{/fapi/v1/markPriceKlines} and
\texttt{/fapi/v1/indexPriceKlines} (note: the index klines endpoint uses
\texttt{pair=BTCUSDT} not \texttt{symbol=BTCUSDT}).

\begin{align}
  b^{\text{prem}}_t &= \frac{M_t - I_t}{I_t} \\[4pt]
  b^z_t &= \frac{b^{\text{prem}}_t - \mu^{(480)}_t}{\sigma^{(480)}_t}
  \quad\text{(8h rolling z-score)} \\[4pt]
  \hat{b}_t &= F_{480}^{-1}(b^{\text{prem}}_t) - 0.5
  \quad\text{(8h rolling percentile rank, centred)}
\end{align}
where $\mu^{(480)}_t$ and $\sigma^{(480)}_t$ are the rolling 480-bar
(8-hour) mean and standard deviation of the raw basis, and
$F_{480}^{-1}(\cdot)$ denotes the empirical CDF within the same window.
Both $b^z_t$ and $\hat{b}_t$ are zero-centred representations of how
extreme the current basis is relative to recent history.

Training period: 2023-01-01 to 2023-12-31 ($n = 525{,}121$ 1-minute bars).
OOS period: 2024-01-01 to 2024-12-31 ($n = 527{,}038$).

\subsection{IC Profile}

\begin{table}[H]
\centering
\caption{IC across horizons for continuous premium basis (training sample)}
\label{tab:e_ic}
\begin{tabular}{lrrrr}
\toprule
Signal & Horizon & IC & $|\IC|/\IC^*_{\text{taker}}$ & $|\IC|/\IC^*_{\text{maker}}$ \\
\midrule
$b^{\text{prem}}$ & 1h  & $-0.024$ & $0.09\times$ & $0.14\times$ \\
$b^{\text{prem}}$ & 4h  & $-0.026$ & $0.19\times$ & $0.31\times$ \\
$b^{\text{prem}}$ & 8h  & $-0.040$ & $0.41\times$ & $0.69\times$ \\
\midrule
$b^z$             & 1h  & $-0.049$ & $0.18\times$ & $0.30\times$ \\
$b^z$             & 4h  & $-0.037$ & $0.27\times$ & $0.44\times$ \\
$b^z$             & 8h  & $-0.026$ & $0.27\times$ & $0.45\times$ \\
\midrule
$\hat{b}$         & 1h  & $-0.048$ & $0.17\times$ & $0.29\times$ \\
$\hat{b}$         & 4h  & $-0.036$ & $0.26\times$ & $0.43\times$ \\
$\hat{b}$         & 8h  & $-0.026$ & $0.26\times$ & $0.44\times$ \\
\bottomrule
\end{tabular}
\end{table}

The IC is uniformly negative (positive basis predicts negative future return),
confirming the direction of the arbitrage mechanism. The best signal is $b^z$
at 1h with IC $= -0.049$.

\subsection{Statistical Tests}

\textbf{Permutation test} ($N = 2000$, 50k-subsampled):
$b^z$ vs 1h return: IC $= -0.047$, $p = 0.0000$ ($-10.8\sigma$ below null).
All horizon / signal combinations are highly significant ($p < 0.001$),
reflecting the large effective sample size.

\textbf{Block bootstrap} (block $= 60$ bars $= 1$ hour, $B = 500$,
50k-contiguous subsample; p-value computed by shifting bootstrap distribution
to null before testing, consistent with percentile CI):\\
$b^z$ vs 1h: IC $= -0.055$, $p = 0.010$, 95\% CI $= [-0.092,\,-0.012]$ — \emph{excludes zero}.\\
$b^z$ vs 8h: IC $= -0.012$, $p = 0.642$, 95\% CI $= [-0.056,\,+0.040]$ — includes zero.

The 1-hour horizon is statistically robust after block-resampling; the 8-hour
horizon is not.

\subsection{Out-of-Sample Stability}

Walk-forward OOS (30-day windows, 10-day step, $n = 34$ windows):
\begin{itemize}
  \item $b^z$ vs 1h: mean IC $= -0.024$, std $= 0.030$, 76\% of windows negative.
  \item $b^z$ vs 8h: mean IC $= -0.002$, std $= 0.041$, 50\% negative (noise).
\end{itemize}
The 1h signal is directionally consistent in OOS.

\subsection{Comparison with Signal H1 and Conclusion}

\begin{table}[H]
\centering
\caption{Continuous (J) vs.\ settlement-bar (H1) basis at 1h horizon}
\label{tab:e_vs_h1}
\begin{tabular}{lrrrr}
\toprule
Version & $n$ (train) & IC (1h) & $|\IC|/\IC^*_{\text{maker}}$ & OOS $\%$ neg \\
\midrule
H1 (discrete, 8h bars)       & $\approx 500$  & $-0.170$ & $1.02\times$ & $70\%$ \\
J (continuous, 1m bars)      & $525{,}121$    & $-0.049$ & $0.30\times$ & $76\%$ \\
\bottomrule
\end{tabular}
\end{table}

The settlement-bar version (H1) is $3.4\times$ stronger in IC despite having
$1000\times$ fewer observations.  The explanation is structural:
\begin{itemize}
  \item At settlement, the funding payment is \emph{mechanically certain} —
        the spread produces an immediate P\&L effect for all open positions.
        The arbitrage is forced.
  \item In continuous time, the arbitrage is only \emph{probabilistic}:
        arbitrageurs may not be active, and the pressure can take time or
        not materialise.
\end{itemize}
Signal J is statistically real (the mechanism is correct) but the
continuous-time version has insufficient IC to cover maker costs at any
tested horizon (best: $0.45\times$ at 8h).


%% ============================================================
\section{Summary and Discussion}
%% ============================================================

\subsection{Signal Rankings}

\begin{table}[H]
\centering
\caption{All signals: summary of findings.
  $p$-val: permutation test.
  n.r.\ = not reported (formal permutation test was not run;
  directional consistency was used instead).
  OOS: fraction of walk-forward windows with IC in the correct direction.
  -- = test not run or result undefined.
  Phase 2 signals are evaluated on OOS 2024+ only (first-pass screen).}
\label{tab:summary}
\begin{tabular}{llrrrrl}
\toprule
Signal & Family & Best IC & $p$-val & OOS & vs $\IC^*$ (maker) & Phase \\
\midrule
C: CVD Divergence         & Order flow & $-0.056$ & $< 0.001$ & $94.9\%$ & $0.33\times$ & 1 \\
E (VPIN regime)           & Order flow & $-0.355$ & $< 0.001$ & --       & (vol only)   & 1 \\
A: Funding curvature      & Sentiment  & $-0.033$ & $0.268$   & Sign flip & $0.38\times$ & 1 \\
B: Informed flow          & Order flow & $-0.026$ & $< 0.001$ & $92.0\%$ & $0.05\times$ & 1 \\
F: VPIN asymmetry         & Order flow & $-0.050$ & $< 0.001$ & $91.7\%$ & $0.30\times$ & 1 \\
H1: Basis (settlement)    & Arbitrage  & $-0.170$ & $0.021$   & $70\%$\textsuperscript{\dag} & $1.02\times$\textsuperscript{\ddag} & 1 \\
J: Basis (continuous 1m)  & Arbitrage  & $-0.049$ & $< 0.001$ & $76.0\%$ & $0.30\times$ & 1 \\
H2: Funding carry         & Carry      & $-0.067$ & n.r.      & $80.0\%$ & $0.40\times$ & 1 \\
I1: Trapped positioning   & Crowding   & $+0.027$ & $0.348$   & --       & --           & 1 \\
I3: Joint extreme         & Crowding   & $-0.040$ & $0.001$   & --       & $0.24\times$ & 1 \\
\midrule
D4: Funding$\times\Delta$OI & OI + carry & $+0.017$ & $0.634$  & -- & $0.34\times$\textsuperscript{$\ddagger$} & 2 (killed) \\
N1: DVOL level              & Vol regime & $+0.024$ & $0.638$\textsuperscript{$\P$} & -- & $0.86\times$\textsuperscript{$\P$}  & 2 (killed)\textsuperscript{$\P$} \\
N2: DVOL 8h slope           & Vol regime & $+0.078$ & $0.136$\textsuperscript{$\P$} & -- & $2.84\times$\textsuperscript{$\P$}  & 2 (killed)\textsuperscript{$\P$} \\
N3: DVOL 30d $z$-score      & Vol regime & $+0.112$ & $0.039$\textsuperscript{$\P$} & 83\%\textsuperscript{$\S$} & $\mathbf{4.11\times}$\textsuperscript{$\P$} & 2 \textbf{survivor}\textsuperscript{$\star$} \\
\bottomrule
\end{tabular}
\par\smallskip\footnotesize
\textsuperscript{\dag}H1 OOS 70\% is within-2023 walk-forward only; calendar OOS shows
IC$\approx$0 in 2024-H1, IC$=-0.044$ in 2024-H2.  Execution feasibility study finds
pre-settlement IC \emph{inverts} (91\% adverse selection at $t-1$\,min); 00:00\,UTC
subset gives IC$=-0.090$ (0.67$\times$ maker) but is still sub-breakeven.
\textsuperscript{\ddag}Training IC is based on only 185 of 1095 bars (sparse \texttt{basis}
column); 1.02$\times$ is not reliable.  Signal is not viable as a standalone strategy.
Phase 2 IC/ratio reported at best OOS horizon after deep-dive (24h for N-signals).
\textsuperscript{$\ddagger$}D4 ratio of 1.16$\times$ in first-pass screen was a methodological
artefact (1m forward-filled bars inflate effective $n$); at true settlement bars
($n=927$) ratio collapses to 0.34$\times$, $p=0.634$.
\textsuperscript{$\S$}N3 OOS\% is fraction of rolling 90-day windows with IC $>0$
(OOS 2024, 276 windows).  Walk-forward extended to 2025--2026: 2 of 4 half-year OOS
periods pass ($p < 0.05$); see Section~\ref{sec:n3_kill}.
\textsuperscript{$\P$}Rigorous daily bootstrap: non-overlapping daily sample (every 1440 bars),
block-bootstrap $B=1000$, block $=14$\,d, OOS 2024.  N1 ($p=0.638$) and N2 ($p=0.136$)
fail; first-pass 1m-overlapping $p\approx 0$ for both is an artefact of
overlap inflation at the 24h horizon.  N1 and N2 are killed after rigorous testing.
\textsuperscript{$\star$}N3 survives but is regime-dependent: IC near zero when DVOL $<51$ ($p=0.691$);
strong when DVOL $\geq 51$ ($p=0.018$, $4.69\times$); see Section~\ref{sec:n3_kill}.
\end{table}

\subsection{The Microstructure Family Problem}

Signals C, B, F, and (indirectly) E\textsuperscript{vpin} all measure variants of the same
phenomenon: \emph{short-horizon taker-side order flow imbalance that
mean-reverts}. They are mathematically related:
\begin{align*}
  \text{CVD}        &= \sum V^B - \sum V^S \text{ (unnormalised, 60-bar)}\\
  \text{OFI}^Q      &= Q^B/Q - 0.5 \text{ (per-bar)} \\
  \text{IF}^{(60)}  &= z^{(60)}(\bar{q}) \cdot \text{OFI}^Q \text{ (size-weighted OFI)}\\
  A_t               &= \EMA_{50}(V^B/V) - \EMA_{50}(V^S/V) \text{ (smoothed OFI)}
\end{align*}
All are variants of the ratio $V^B/V$ integrated over different time windows
and with different weightings. Their OOS ICs are similar ($-0.020$ to $-0.056$)
because they measure the same underlying quantity. Combining them produces
no material improvement in predictive power.

\subsection{The Research Funnel}

Following the systematic research framework:
\begin{enumerate}
  \item \textbf{Economic mechanism} $\to$ explains \emph{why} the signal should work
  \item \textbf{Measurable signal construction} $\to$ implementable formula
  \item \textbf{Statistical predictability test} $\to$ IC, permutation, bootstrap
  \item \textbf{New signal check} $\to$ not just a repackaged existing signal
  \item \textbf{Breakeven IC estimate before full backtest}
  \item \textbf{Out-of-sample validation}
  \item \textbf{Execution cost adjustment}
  \item \textbf{Stress test by regime}
  \item \textbf{Portfolio combination}
\end{enumerate}

The order-flow signals (C, B, F) fail at steps 4 and 5: they are related to
known microstructure signals and are economically sub-breakeven. Signal A
fails at step 3. Signal H1 clears step 5 in training but fails step 7:
pre-settlement IC inverts (91\% adverse selection at $t-1$\,min), making
maker-only entry structurally impossible. Phase 2 candidates D4 and N1--N3
are the first to clear steps 1--5 with ratios above $1\times$; they now
proceed to steps 6--8 (full OOS walk-forward, bootstrap, execution design).

\subsection{Key Lessons}

\begin{enumerate}
  \item \textbf{Statistical vs economic significance are distinct.}
        A signal with $p < 0.001$ and $94.9\%$ OOS directional stability
        (Signal C) may still lose money. Statistical significance is necessary
        but not sufficient.

  \item \textbf{The microstructure mean-reversion family is over-crowded.}
        BTCUSDT 1-minute order flow is one of the most competitive signals
        in crypto. Market makers with microsecond latency already exploit
        these patterns; taker-cost strategies face an adverse selection
        problem.

  \item \textbf{Look-ahead bias is subtle.}
        Using $r_t^{(h)}$ shifted by one bar as a ``past'' return inflates
        IC from $0.027$ to $0.710$ when $h = 480$ bars. All lookback windows
        must use genuinely historical quantities.

  \item \textbf{Mechanism-driven signals differ from pattern signals.}
        The basis dislocation signal (H1) is driven by a deterministic
        arbitrage force at a scheduled time (funding settlement). It does not
        require the pattern to persist; it requires the mechanism to persist,
        which is more likely as long as the funding system exists.

  \item \textbf{Maker execution changes the economics substantially.}
        At 1h horizon, switching from taker to maker reduces the breakeven IC
        from $0.279$ to $0.168$, a $40\%$ reduction. Signal H1 crosses
        breakeven only at maker costs.

  \item \textbf{Options market implied volatility encodes regime information.}
        The Deribit DVOL index --- a model-free 30-day IV measure --- is the
        strongest survivor of the programme, with N3 ($z$-score variant)
        achieving $4.11\times$ maker breakeven at 24h on non-overlapping daily
        returns ($p = 0.039$). The empirical IC is \emph{positive}: when DVOL
        is high relative to its recent history, BTC perpetual returns over the
        next 24h are systematically above average. This is consistent with a
        \emph{fear-resolution} effect --- elevated IV signals market stress;
        the subsequent normalisation as fear subsides produces positive returns
        for those who buy into panic. It is not a negative-carry mechanism.
\end{enumerate}

\subsection{Phase 2 Quick Screen}
\label{sec:phase2screen}

\textbf{Phase 2 ran a first-pass IC screen on nine signals across five mechanism
families.}  Pass criterion: $|\IC|/\IC^*_{\text{maker}} > 0.5$ \emph{and}
permutation $p < 0.10$ at any tested horizon.  OOS period: 2024-01-01 onward
(527\,040 one-minute bars).

\begin{table}[ht]
\centering
\caption{Phase 2 first-pass IC screen results (OOS 2024--2026)}
\label{tab:phase2screen}
\begin{tabular}{llrrrl}
\toprule
ID & Signal & Best Hz & $|\IC|/\IC^*$ & $p$ & Decision \\
\midrule
N3 & DVOL 30-day $z$-score                       & 24h & 3.887 & 0.000 & \textbf{pass} \\
N2 & DVOL 8h slope                               & 24h & 2.371 & 0.000 & \textbf{pass} \\
N1 & DVOL level (raw)                            & 24h & 1.197 & 0.000 & \textbf{pass} \\
\midrule
D4 & Funding $\times$ $\Delta$OI at settlement   & 4h  & 0.339 & 0.634 & kill (artefact)\textsuperscript{$\ddagger$} \\
L1 & Binance$-$Bybit $\Delta$funding             & 8h  & 0.687 & 0.560 & kill ($p$ fails) \\
D1 & $\Delta$OI velocity ($z$-scored)            & 4h  & 0.125 & 0.064 & kill (ratio fails) \\
D2 & $\Delta^2$OI acceleration                   & 4h  & 0.017 & 0.618 & kill \\
K  & Liq.\ imbalance proxy (klines)              & 4h  & 0.073 & 0.148 & kill \\
L2 & Binance$-$OKX $\Delta$funding               & ---  & ---   & ---   & kill (OKX API $\leq$90d) \\
\bottomrule
\end{tabular}
\end{table}

\noindent \textbf{Data notes.}
Binance liquidation snapshot data is no longer available on
\texttt{data.binance.vision}; Signal~K was evaluated using a volume-imbalance
proxy from 1-minute klines.  The OKX public funding-rate API retains only
$\approx$90 days of history, making a multi-year OOS test impossible.
Bybit funding (L1) provides 411 aligned settlement bars in OOS but the
permutation $p = 0.56$ indicates the observed IC is consistent with noise.

\textbf{First-pass survivors}: D4 and the full DVOL family (N1, N2, N3)
cleared the hurdle.  After deep-dive validation, \textbf{D4 was killed}: its
1.16$\times$ ratio was an artefact of forward-filling settlement values into
1-minute bars (inflating effective $n$ without adding independent information).
At true settlement bars ($n = 927$) D4 collapses to $0.34\times$, $p = 0.634$.
\textbf{The sole Phase 2 survivor is the DVOL family (N1, N2, N3)}, subject to
the caveat that the 24h IC ratios in Table~\ref{tab:phase2screen} use overlapping
1-minute returns and may overstate significance --- see Section~\ref{sec:sigN}
for non-overlapping validation.

\subsection{Phase 2 Deep Dive: DVOL Volatility Regime (Signal N)}
\label{sec:sigN}

\subsubsection{Mechanism}

The Deribit BTCVOL index (\emph{DVOL}) is a 30-day model-free implied
volatility measure computed from options prices, analogous to VIX for equity.
Three distinct regimes drive perpetual returns:

\begin{enumerate}
  \item \textbf{DVOL level (N1)}: elevated realised IV is associated with
        negative carry environments and forced de-leveraging.  The perpetual
        trades at a structural discount to spot (funding inverts) when dealers
        hedge short-gamma books, depressing mark price.

  \item \textbf{DVOL slope (N2)}: a rising DVOL (front IV increasing faster
        than the term structure) signals escalating near-term fear; a falling
        DVOL signals the \emph{release} of that fear (vol-selling rally).
        The 8-hour slope is matched to the funding-settlement cycle.

  \item \textbf{DVOL $z$-score (N3)}: the 30-day rolling $z$-score separates
        the level effect from the trend, capturing how unusual the current
        vol environment is relative to recent history.  Empirically the IC
        is \emph{positive} (high $z$-score $\Rightarrow$ positive future
        return) --- consistent with a fear-resolution effect rather than
        negative carry, as discussed in Section~\ref{sec:sigN_deepdive}.
\end{enumerate}

\subsubsection{Quick-Screen Results}

All three DVOL variants clear the $0.5\times$ ratio and $p < 0.10$ hurdles
at the 8-hour horizon (Table~\ref{tab:phase2screen}).  N3 achieves
$|\IC|/\IC^* = 1.62$ with $p \approx 0$, indicating the signal carries
economically meaningful information net of round-trip maker costs.

\subsubsection{Rigorous Validation: Non-Overlapping Returns}
\label{sec:sigN_deepdive}

The initial deep-dive in Section~\ref{sec:sigN} used all 527\,040 OOS
1-minute bars as independent observations for the 24h forward return.
Adjacent observations share $1439/1440 = 99.9\%$ of their 24h return
window, reducing the effective sample size from 527k to approximately
366 independent days.

\paragraph{Overlap inflation check.}
Sampling every 1440 bars (one observation per day, $n = 365$) and comparing
to the full-overlap series:
\begin{center}
\begin{tabular}{lrrr}
\toprule
Series & $n$ & IC (24h) & Ratio \\
\midrule
1m-overlapping (all bars)      & 525,600 & $+0.109$ & $3.89\times$ \\
Non-overlapping (daily sample) &     365 & $+0.112$ & $4.11\times$ \\
\bottomrule
\end{tabular}
\end{center}
The inflation factor is $0.95\times$ --- the overlapping series gives a
\emph{lower} IC than the non-overlapping series, not higher.  There is no
evidence of overlap inflation in this signal.

\paragraph{Block bootstrap on daily series.}
Table~\ref{tab:n3_rigorous} reports block-bootstrap results using
non-overlapping daily observations (block size $= 14$ days, $B = 1000$).

\begin{table}[ht]
\centering
\caption{Rigorous block-bootstrap on non-overlapping daily observations (OOS 2024, $B=1000$, block=14d)}
\label{tab:n3_rigorous}
\begin{tabular}{llrrrrl}
\toprule
Signal & Hz & IC & $p$ & 95\% CI & Ratio & Verdict \\
\midrule
N3 ($z$-score) & 24h & $+0.112$ & $\mathbf{0.039}$ & $[-0.103, +0.110]$ & $4.11\times$ & \textbf{pass} \\
N3 ($z$-score) & 48h & $+0.158$ & $\mathbf{0.002}$ & $[-0.101, +0.106]$ & $7.99\times$ & \textbf{pass} \\
N2 (slope)     & 24h & $+0.078$ & $0.136$           & $[-0.109, +0.105]$ & $2.84\times$ & fail ($p$) \\
N2 (slope)     & 48h & $+0.050$ & $0.351$           & $[-0.101, +0.096]$ & $2.53\times$ & fail ($p$) \\
N1 (level)     & 24h & $+0.024$ & $0.638$           & $[-0.101, +0.114]$ & $0.86\times$ & fail \\
N1 (level)     & 48h & $+0.039$ & $0.454$           & $[-0.096, +0.106]$ & $1.97\times$ & fail ($p$) \\
\bottomrule
\end{tabular}
\end{table}

\noindent \textbf{Only N3 passes rigorous testing.}  N1 and N2 are not
statistically significant at the daily granularity after controlling for
temporal autocorrelation.  N3 is significant at both 24h ($p = 0.039$) and
48h ($p = 0.002$).

\paragraph{N3 incremental over N1/N2.}
N1--N3 Spearman correlations: $\rho(N1, N3) = 0.56$; $\rho(N2, N3) = 0.21$;
$\rho(N1, N2) = 0.07$.  After orthogonalising N3 against N1, the residual IC
remains $+0.112$ ($4.09\times$) --- N1 carries virtually no information beyond
what N3 already captures.  \textbf{N3 is the only signal that needs to be
traded.}

\subsubsection{Walk-Forward, Total PnL, and Long-Short Performance}

\paragraph{IC sign.}
The IC for N3 is \emph{positive} (high $z$-score $\Rightarrow$ positive future
return) --- consistent with a \emph{fear-resolution} mechanism.  When DVOL is
well above its 30-day mean, the market is near a local fear peak; as the vol
regime normalises, BTC recovers.

\paragraph{Walk-forward IC by half-year (daily non-overlapping).}
Training (2023) yields IC $= +0.005$ ($0.14\times$) --- near zero, confirming
no in-sample overfitting.  OOS 2024-H1: IC $= +0.177$ ($6.65\times$).
OOS 2024-H2: IC $= +0.091$ ($3.24\times$).  Both OOS halves are strong,
consistent, and substantially above the maker breakeven.  OFI data currently
ends 2024-12-31; 2025--2026 validation requires extending the dataset.

\paragraph{Total perpetual PnL.}
Including funding payments received/paid by a long position over 24h,
the IC barely changes: $+0.112 \to +0.110$ ($4.02\times$).  The mean daily
funding cost for a long in 2024 was $-3.3$\,bp (negative: shorts paid longs
on average), so funding slightly \emph{improved} total PnL relative to pure
price return.

\paragraph{Net-of-cost quintile returns (daily, total PnL).}
Table~\ref{tab:n3_quintile} reports N3 quintile returns at the 24h horizon
for total PnL (price + funding), net of maker round-trip cost.

\begin{table}[ht]
\centering
\caption{N3 quintile total PnL at 24h horizon, net of maker cost (OOS 2024, $n=365$ days)}
\label{tab:n3_quintile}
\begin{tabular}{lrrrr}
\toprule
Quintile & $n$ & Mean N3$z$ & Gross 24h & Net (maker) \\
\midrule
Q1 (lowest)  & 73 & $-1.76$ & $-17.7$\,bp & $-23.7$\,bp \\
Q2           & 73 & $-0.84$ & $-12.8$\,bp & $-18.8$\,bp \\
Q3           & 73 & $-0.07$ & $+1.9$\,bp  & $-4.1$\,bp  \\
Q4           & 73 & $+0.78$ & $+41.7$\,bp & $+35.7$\,bp \\
Q5 (highest) & 73 & $+2.00$ & $+78.0$\,bp & $+72.0$\,bp \\
\midrule
Q5$-$Q1 spread & & & $+95.7$\,bp & $+83.7$\,bp \\
\bottomrule
\end{tabular}
\end{table}

\paragraph{Long-short rule performance (OOS 2024).}
A simple rule --- long when $N3z > \theta$, short when $N3z < -\theta$,
flat otherwise --- using daily rebalancing and maker-cost round-trips:

\begin{center}
\begin{tabular}{lrrrrrr}
\toprule
$\theta$ & Active days & Trades & Net PnL & Sharpe & Max DD \\
\midrule
0.00 & 365 & 42  & $+8{,}591$\,bp & 1.68 & $2{,}498$\,bp \\
0.25 & 323 & 73  & $+9{,}358$\,bp & 1.97 & $2{,}064$\,bp \\
0.50 & 272 & 96  & $+10{,}032$\,bp & 2.24 & $2{,}064$\,bp \\
0.75 & 228 & 82  & $+9{,}305$\,bp & \textbf{2.37} & $1{,}924$\,bp \\
1.00 & 188 & 82  & $+7{,}776$\,bp & 2.26 & $1{,}849$\,bp \\
\bottomrule
\end{tabular}
\end{center}

\noindent Sharpe peaks at $\theta = 0.75$ ($\text{Sharpe} = 2.37$, net PnL
$= 9{,}305$\,bp over 2024, max drawdown $= 1{,}924$\,bp).  All thresholds
are profitable net of maker costs.

\paragraph{Event robustness.}
Excluding the 6 event days where $|r_{24h}| > 3\,\sigma$ ($> 827$\,bp),
the IC drops modestly from $+0.112$ to $+0.103$ ($3.44\times$).  The signal
does not depend on a handful of extreme outlier events.

\paragraph{Direction vs magnitude.}
N3 predicts \emph{direction} (IC $= +0.112$, ratio $= 4.11\times$) more
strongly than \emph{magnitude} (IC($N3z, |r_{24h}|) = +0.087$, $2.18\times$),
confirming it is a directional signal rather than purely a volatility-timing
tool.

\subsubsection{Verdict on Signal N}

\textbf{N3 is the strongest signal in the programme and survives rigorous
validation, with important regime caveats discovered in the 2025--2026
kill-attempt (Section~\ref{sec:n3_kill}).}

Summary of validated results (OOS 2024, $n=365$ daily non-overlapping):
\begin{itemize}[noitemsep]
  \item No overlap inflation: non-overlapping IC ($4.11\times$) matches or
        exceeds overlapping IC ($3.89\times$).
  \item Block-bootstrap on daily series ($B=1000$, block $=14$\,d):
        $p = 0.039$ at 24h, $p = 0.002$ at 48h.
  \item Walk-forward: near-zero training IC ($0.14\times$), strong OOS
        ($6.65\times$ in 2024-H1, $3.24\times$ in 2024-H2).
  \item Total PnL including funding: IC ratio $4.02\times$ (funding helps).
  \item Long-short rule: Sharpe $= 2.37$ ($\theta = 0.75$), net PnL
        $+9{,}305$\,bp in 2024, max drawdown $1{,}924$\,bp.
  \item Event-robust: IC $= 3.44\times$ ex-event days.
  \item N1 and N2 do \emph{not} survive rigorous daily bootstrap; N3 carries
        essentially all the information.
\end{itemize}

Additional findings from 2025--2026 extension:
\begin{itemize}[noitemsep]
  \item 2025-H1 passes independently ($p = 0.031$, IC $= 5.32\times$) ---
        signal is not a single-year fluke.
  \item 2025-H2 and 2026-YTD fail; signal is intermittent by period.
  \item Regime conditioning is essential: IC near zero in low-DVOL
        environments ($< 51$, $p = 0.691$) but strong in high-DVOL
        ($\geq 51$, $p = 0.018$, $4.69\times$).
  \item Edge is distributed, not concentrated in outlier days.
\end{itemize}

\subsection{Phase 2 Deep Dive: OI--Funding Combo (Signal D4)}
\label{sec:sigD4}

\subsubsection{Mechanism}

D4 combines two positioning signals at each 8-hour funding settlement:
\[
  D4_t = F_t \times z(\Delta\text{OI}_{8h,t})
\]
where $F_t$ is the contemporaneous funding rate and $z(\Delta\text{OI})$ is
the standardised change in open interest since the previous settlement.
The interaction captures the joint state of \emph{how much the market is
paying for leverage} and \emph{whether new leverage is being added or
removed}.  A high funding rate coinciding with rising OI signals crowded
long positioning; the perpetual tends to soften once the settlement resolves
the carry.

\subsubsection{Quick-Screen and Deep-Dive Results}

D4 achieved $|\IC|/\IC^* = 1.16$ at 4h in the first-pass screen, but this
was computed on all 527\,040 OOS 1-minute bars rather than on the
settlement bars only.  When evaluated correctly --- at the $N=927$
Binance 8-hour settlement timestamps in OOS --- the signal collapses:
$\IC = +0.017$, ratio $= 0.34\times$, block-bootstrap $p = 0.634$.
The apparent 1.16$\times$ in the quick screen was a methodological artefact:
the 1m-forward-filled OI and funding values repeat the same signal value
over hundreds of bars before the next settlement, inflating the effective
sample size while not adding independent observations.

\textbf{D4 is killed.}  The interaction of OI change and funding rate carries
no significant predictive power at actual settlement events with 927 OOS
observations.  This is consistent with D1 and D2 also failing: OI dynamics
do not appear to predict BTCUSDT perpetual returns in the 2024 OOS period.

%% ============================================================
\subsection{N3 Kill-Attempt: 2025--2026 Out-of-Sample}
\label{sec:n3_kill}
%% ============================================================

Kline and funding data were extended through 2026-05-13 (1{,}768{,}728 bars
from 2023-01-01).  The same non-overlapping daily protocol from
Section~\ref{sec:sigN_deepdive} was applied.  The objective was to break N3
if it is spurious.

\subsubsection{Walk-Forward IC Across All Periods}

\begin{table}[ht]
\centering
\caption{N3 walk-forward IC: non-overlapping daily, block-bootstrap $p$ (block=14 days)}
\label{tab:n3_walkforward}
\begin{tabular}{lrrrrl}
\toprule
Period & $n$ & IC & $p$ & Ratio & Verdict \\
\midrule
Train 2023     & 365 & $+0.005$ & 0.933 & $0.14\times$ & no overfit \\
OOS 2024-H1    & 182 & $+0.177$ & 0.018 & $6.65\times$ & \textbf{PASS} \\
OOS 2024-H2    & 184 & $+0.093$ & 0.224 & $3.31\times$ & weak \\
OOS 2025-H1    & 181 & $+0.164$ & 0.031 & $5.32\times$ & \textbf{PASS} \\
OOS 2025-H2    & 184 & $+0.017$ & 0.813 & $0.43\times$ & fail \\
OOS 2026-YTD   & 132 & $-0.068$ & 0.441 & $2.45\times$ & fail \\
\bottomrule
\end{tabular}
\end{table}

The first headline is that N3 passes two independent H1 windows
(2024-H1 and 2025-H1) with similar IC ($\approx 0.17$) and $p < 0.05$.
Training IC is near zero ($0.14\times$), confirming no in-sample overfit.
However, both H2 windows and the 2026 YTD fail: the signal is intermittent,
not consistently positive across all periods.

\subsubsection{Year-by-Year Backtest}

\begin{table}[ht]
\centering
\caption{Year-by-year backtest: long-short rule $\theta=0.75$, maker cost, daily total PnL}
\label{tab:n3_yearly}
\begin{tabular}{lrrrrr}
\toprule
Year & $n$ & Sharpe & PnL (bp) & MaxDD (bp) & $n_{\text{trades}}$ \\
\midrule
2023 (train) & 365 & $+0.41$ & $+1{,}650$ & $-1{,}962$ & 77 \\
2024 (OOS)   & 366 & $+1.88$ & $+9{,}375$ & $-2{,}050$ & 84 \\
2025 (OOS)   & 365 & $+0.20$ & $+784$     & $-4{,}638$ & 89 \\
2026 (OOS)   & 132 & $-0.29$ & $-563$     & $-2{,}230$ & 23 \\
\bottomrule
\end{tabular}
\end{table}

2024 is strong (Sharpe 1.88); 2025 is marginal (+0.20); 2026 YTD is slightly
negative.  The signal is not dead but it has clearly weakened beyond 2024.

\subsubsection{Regime Conditioning: DVOL Level}

The most important finding from the kill-attempt is that the IC is almost
entirely confined to \emph{high-DVOL} environments.

\begin{table}[ht]
\centering
\caption{N3 IC by DVOL regime (OOS 2024+, non-overlapping daily)}
\label{tab:n3_regime}
\begin{tabular}{lrrrr}
\toprule
Regime & $n$ & IC & Ratio & $p$ \\
\midrule
Low DVOL ($< 51$, median)      & 430 & $+0.019$ & $0.51\times$ & 0.691 \\
High DVOL ($\geq 51$, median)  & 433 & $+0.119$ & $4.69\times$ & 0.018 \\
High DVOL ($\geq 57$, 75th)    & 216 & $+0.153$ & $6.47\times$ & 0.031 \\
\bottomrule
\end{tabular}
\end{table}

In low-DVOL environments (DVOL $< 51$) the IC is statistically and economically
zero ($p = 0.691$).  When DVOL is at or above its median (51), the signal
recovers strongly ($p = 0.018$, $4.69\times$ maker breakeven).  At the 75th
percentile (DVOL $\geq 57$) the IC reaches $0.153$ ($p = 0.031$).

This is the primary explanation for the period-to-period variability: the H2
and 2026 periods likely contain more low-DVOL bars (suppressed volatility
during a risk-on environment), muting the signal.

\subsubsection{Concentration and Event Robustness}

Using the $\theta = 0.75$ long-short rule on OOS 2024+ data:
\begin{itemize}[noitemsep]
  \item Top-5 days \emph{drag} performance by $-1{,}185$\,bp; removing them
        \emph{increases} Sharpe from $+0.88$ to $+1.10$.  The edge is
        distributed, not concentrated in outlier days.
  \item Ex-tail-events (excluding top-5\% by $|r_{24h}|$): IC improves from
        $2.50\times$ to $2.11\times$ with $p = 0.022$.  Signal not dependent
        on extreme-event days.
\end{itemize}

\subsubsection{Threshold Sensitivity (OOS 2024+)}

\begin{table}[ht]
\centering
\caption{Long-short backtest threshold sensitivity (OOS 2024+, total PnL, maker cost)}
\label{tab:n3_thresh2}
\begin{tabular}{rrrrrr}
\toprule
$\theta$ & $n_{\text{traded}}$ & Sharpe & PnL (bp) & MaxDD (bp) \\
\midrule
0.00 & 863 & $+0.59$ & $+8{,}103$  & $-5{,}564$ \\
0.50 & 644 & $+0.81$ & $+9{,}867$  & $-5{,}735$ \\
0.75 & 544 & $+0.88$ & $+9{,}595$  & $-5{,}055$ \\
1.00 & 455 & $+0.93$ & $+9{,}209$  & $-4{,}049$ \\
1.25 & 353 & $+0.96$ & $+8{,}320$  & $-4{,}086$ \\
1.50 & 239 & $+0.54$ & $+4{,}023$  & $-3{,}928$ \\
\bottomrule
\end{tabular}
\end{table}

Sharpe peaks at $\theta = 1.25$ ($+0.96$, 353 active days out of 863).
Tighter thresholds reduce drawdown while approximately maintaining PnL
until $\theta = 1.50$ where sample size collapses.

\subsubsection{Kill-Attempt Verdict}

\textbf{N3 is not killed, but it is clearly regime-dependent.}

\begin{itemize}[noitemsep]
  \item \textbf{Survives} two independent H1 windows (2024-H1 and 2025-H1)
        with $p < 0.05$ each --- not a single-year fluke.
  \item \textbf{Weakens} materially in H2 periods and 2026 YTD --- signal is
        not uniformly persistent.
  \item \textbf{Regime filter required}: IC is statistically zero in low-DVOL
        environments (DVOL $< 51$, $p = 0.691$) and strong in high-DVOL
        ($p = 0.018$).  A conditional strategy that only trades when
        DVOL $\geq$ its rolling median is well-motivated.
  \item \textbf{Not concentrated}: edge is distributed across many days;
        top-5 outlier days hurt rather than help.
  \item \textbf{Total OOS 2024+ Sharpe}: $+0.88$ at $\theta = 0.75$,
        $+0.93$ at $\theta = 1.00$ --- modest but positive with 863 daily
        observations.  The DVOL-conditioned strategy (Section~\ref{sec:n3_regime_filter}) materially improves this.
\end{itemize}

The primary open question is whether the H2/2026 weakness reflects low-DVOL
regimes (regime explanation) or genuine signal decay (structural explanation).
This is answered directly in Section~\ref{sec:n3_regime_filter} below.

\subsubsection{DVOL-Conditioned Strategy}
\label{sec:n3_regime_filter}

\paragraph{Why the weak periods failed.}
The DVOL level distribution by period reveals the mechanism directly:

\begin{table}[ht]
\centering
\caption{DVOL level during each period — proportion of high-IV days}
\label{tab:dvol_dist}
\begin{tabular}{lrrrr}
\toprule
Period & $n$ & DVOL mean & DVOL median & Pct $\geq 51$ \\
\midrule
Train 2023   & 365 & 49.5 & 50.5 & 47\% \\
OOS 2024-H1  & 182 & 59.3 & 56.1 & 74\% \\
OOS 2024-H2  & 184 & 56.1 & 56.5 & 88\% \\
OOS 2025-H1  & 181 & 50.5 & 49.8 & 47\% \\
OOS 2025-H2  & 184 & \textbf{41.6} & \textbf{39.5} & \textbf{4\%} \\
OOS 2026-YTD & 133 & 47.5 & 46.1 & 42\% \\
\bottomrule
\end{tabular}
\end{table}

2025-H2 is the most extreme case: mean DVOL $= 41.6$, median $= 39.5$, and
only 4\% of days had DVOL $\geq 51$.  The BTC bull market of mid-2025
suppressed implied volatility, eliminating the fear-resolution mechanism
entirely.  There was no alpha to capture because the regime that produces it
did not exist.  This is a regime explanation, not signal decay.

\paragraph{DVOL threshold sweep (OOS 2024+).}

\begin{table}[ht]
\centering
\caption{DVOL entry filter sweep: Sharpe and PnL on active days (OOS 2024+, $n=863$ days, $\theta=0.75$, maker cost)}
\label{tab:dvol_sweep}
\begin{tabular}{rrrrrr}
\toprule
DVOL filter & $n_{\text{active}}$ & IC & $p$ & Sharpe & PnL (bp) \\
\midrule
None        & 863 & $+0.072$ & 0.033 & $+0.88$ & $+9{,}595$  \\
$\geq 46$   & 590 & $+0.122$ & 0.004 & $+1.71$ & $+13{,}907$ \\
$\geq 51$   & 445 & $+0.124$ & 0.006 & $+1.68$ & $+10{,}955$ \\
$\geq 54$   & 333 & $+0.117$ & 0.035 & $+2.20$ & $+10{,}895$ \\
$\geq 57$   & 210 & $+0.146$ & 0.035 & $+\mathbf{3.46}$ & $+12{,}246$ \\
$\geq 60$   & 121 & $+0.150$ & 0.120 & $+3.32$ & $+7{,}472$  \\
\bottomrule
\end{tabular}
\end{table}

Sharpe increases monotonically as the DVOL threshold rises, peaking near
DVOL $\geq 57$--60 before sample size becomes restrictive.  The sweet spot
is DVOL $\geq 54$--57: Sharpe $= 2.20$--$3.46$, with maximum drawdown
shrinking from $-5{,}055$\,bp (unfiltered) to $-1{,}624$\,bp.

\paragraph{Per-period conditioned strategy (DVOL $\geq 51$).}

\begin{table}[ht]
\centering
\caption{Long-short Sharpe unconditional vs DVOL$\geq$51 by period ($\theta=0.75$, maker cost)}
\label{tab:dvol_cond}
\begin{tabular}{lrrrr}
\toprule
Period & Sharpe (all) & Sharpe ($\geq$51) & PnL (bp, filtered) & Days passing \\
\midrule
Train 2023   & $+0.41$ & $+1.85$ & $+3{,}812$ & 47\% \\
OOS 2024-H1  & $+2.26$ & $+2.32$ & $+4{,}932$ & 74\% \\
OOS 2024-H2  & $+1.46$ & $+1.51$ & $+3{,}018$ & 88\% \\
OOS 2025-H1  & $+1.10$ & $+1.87$ & $+2{,}187$ & 47\% \\
OOS 2025-H2  & $-1.03$ & $-0.43$ & $-53$       & \textbf{4\%} \\
OOS 2026-YTD & $-0.29$ & $+0.84$ & $+877$      & 42\% \\
\bottomrule
\end{tabular}
\end{table}

The DVOL filter correctly identifies that 2025-H2 is nearly off-limits (4\%
filter pass rate, strategy mostly in cash), and recovers 2026-YTD from
Sharpe $-0.29$ to $+0.84$.

Annual summary at DVOL $\geq 51$:
\begin{itemize}[noitemsep]
  \item \textbf{2024}: Sharpe $+1.89 \to +1.92$; 81\% of days pass filter
        (2024 was already a high-IV year --- filter barely changes anything).
  \item \textbf{2025}: Sharpe $+0.20 \to +1.65$; only 25\% of days pass
        filter but those days deliver strong edge.
  \item \textbf{2026 YTD}: Sharpe $-0.29 \to +0.84$; 42\% pass filter.
\end{itemize}

\paragraph{Regime filter verdict.}

\textbf{The regime explanation is confirmed.  The H2/2026 weakness is entirely
explained by low-DVOL environments; it is not signal decay.}

Specifically:
\begin{itemize}[noitemsep]
  \item Signal mechanism only fires in fear regimes (elevated IV): when
        DVOL is suppressed by a bull market, the fear-resolution return
        pattern simply does not exist.
  \item DVOL $\geq 51$ is a clean regime filter: pass rate is 74--88\% in
        strong periods and $\leq 4\%$ during the 2025 bull suppression.
        The filter is behaving exactly as expected.
  \item DVOL $\geq 54$ appears to be the practical sweet spot: Sharpe $+2.20$,
        $n = 333$ active OOS days (still large enough to be statistically
        meaningful), max drawdown reduced by $68\%$ vs unfiltered.
  \item The rolling-median filter (DVOL $\geq$ 30-day median) is \emph{inferior}
        to the fixed-threshold approach because it allows trading in
        low-IV regimes when the local median happens to be low.
        A fixed absolute threshold (51 or 54) is more robust.
\end{itemize}

%% ============================================================
\subsection{Position-Level Strategy Backtest}
\label{sec:backtest}
%% ============================================================

The frozen rule---N3$z > 0.75$ AND DVOL $\geq 54$, hold 24\,h, maker
execution---was stress-tested with full per-trade PnL accounting:

\begin{equation*}
  \text{Net return}_t = \text{pos}_t \times
    \bigl(r^{24\text{h}}_t - f^{24\text{h}}_t\bigr) - c_{\text{entry}}
\end{equation*}

where $f^{24\text{h}}_t$ is the sum of the three 8-hour funding settlements
in the next 24\,h (paid by longs, received by shorts) and $c_{\text{entry}}$
is the one-way cost charged only when a new position is opened.
Maker round-trip cost: 6\,bp.

\subsubsection*{Master Trade Log (OOS 2024--2026)}

Over 863 OOS calendar days the frozen rule triggered 199 trades
(16.2\% exposure):

\begin{center}
\begin{tabular}{lr}
\hline
Metric & Value \\
\hline
Trades & 199 \\
Sharpe (annualised) & $+2.95$ \\
Net PnL & $+11{,}228$\,bp \\
Max drawdown & $-1{,}612$\,bp \\
Win rate & 53.8\% \\
Avg winning trade & $+262$\,bp \\
Avg losing trade & $-182$\,bp \\
Exposure & 16.2\% of days \\
\hline
\end{tabular}
\end{center}

\subsubsection*{Period Breakdown}

\begin{table}[H]
\centering
\caption{Position-level backtest by period (frozen rule, maker cost).}
\label{tab:backtest_periods}
\begin{tabular}{lrrrrrr}
\hline
Period & $n$ & Sharpe & PnL (bp) & MaxDD (bp) & Win\% & Exp\% \\
\hline
Train 2023   &  77 & $+2.89$ & $+3{,}641$ & $-1{,}154$ & 51.9\% & 6.3\% \\
OOS 2024-H1  &  81 & $+2.57$ & $+4{,}071$ & $-1{,}496$ & 51.9\% & 6.6\% \\
OOS 2024-H2  &  63 & $+3.01$ & $+3{,}259$ & $-1{,}149$ & 50.8\% & 5.1\% \\
OOS 2025-H1  &  38 & $+2.31$ & $+1{,}769$ & $-1{,}346$ & 65.8\% & 3.1\% \\
OOS 2025-H2  &   4 & $+3.41$ &   $+171$   &    $-55$   & 50.0\% & 0.3\% \\
OOS 2026-YTD &  13 & $+5.93$ & $+1{,}958$ &   $-605$   & 46.2\% & 1.1\% \\
\hline
\end{tabular}
\end{table}

Every period shows a positive Sharpe.  The 2025-H2 entry count falls to 4
trades because DVOL was suppressed (mean = 41.6); the regime filter
correctly kept the strategy flat.  Year-by-year totals are stable:
$+2.89$ (2023), $+2.75$ (2024), $+2.38$ (2025), $+5.93$ (2026,
$n = 13$, interpret with caution).

\subsubsection*{Event Dependence}

\begin{table}[H]
\centering
\caption{Event robustness (OOS 2024+, frozen rule, maker cost).}
\label{tab:backtest_events}
\begin{tabular}{lrrrr}
\hline
Sample & $n$ & Sharpe & PnL (bp) & MaxDD (bp) \\
\hline
All OOS days               & 199 & $+2.95$ & $+11{,}228$ & $-1{,}612$ \\
Ex top/bottom 1\% $|r|$   & 194 & $+2.50$ & $+8{,}133$  & $-1{,}659$ \\
Ex 8 known macro events    & 193 & $+2.98$ & $+10{,}767$ & $-1{,}612$ \\
\hline
\end{tabular}
\end{table}

Removing the eight largest-event days (ETF approval, halving, election,
tariff crash, etc.) leaves Sharpe $+2.98$---the strategy does not depend
on these events.  Removing the top/bottom 1\% return days reduces Sharpe
to $+2.50$, indicating that some edge is earned on large-return days, but
the strategy survives without them.

\subsubsection*{Threshold Stability}

\begin{table}[H]
\centering
\caption{Threshold sensitivity grid (OOS 2024+, maker cost). Frozen
         parameters marked $\star$.}
\label{tab:backtest_thresh}
\begin{tabular}{ccrrrr}
\hline
$\theta_{N3z}$ & DVOL$_{\min}$ & $n$ & Sharpe & PnL (bp) & MaxDD (bp) \\
\hline
0.50 & 51 & 306 & $+2.13$ & $+12{,}752$ & $-2{,}074$ \\
0.75 & 51 & 252 & $+2.34$ & $+11{,}444$ & $-2{,}302$ \\
1.00 & 51 & 205 & $+2.64$ & $+10{,}257$ & $-1{,}935$ \\
\hline
0.50 & 54 & 238 & $+2.37$ & $+11{,}076$ & $-2{,}326$ \\
$\mathbf{0.75}$ & $\mathbf{54}$ & $\mathbf{199}$ & $\mathbf{+2.95}$ & $\mathbf{+11{,}228}$ & $\mathbf{-1{,}612}$ $\star$ \\
1.00 & 54 & 168 & $+3.46$ & $+10{,}561$ & $-1{,}425$ \\
\hline
0.50 & 57 & 169 & $+3.67$ & $+12{,}937$ & $-1{,}680$ \\
0.75 & 57 & 144 & $+4.27$ & $+12{,}429$ & $-1{,}099$ \\
1.00 & 57 & 122 & $+4.39$ & $+10{,}482$ &   $-961$   \\
\hline
\end{tabular}
\end{table}

All nine cells show positive Sharpe (range $+2.13$ to $+4.39$).  The frozen
rule sits in the middle of the table; it is not cherry-picked.  Higher
DVOL thresholds yield higher Sharpe at the cost of fewer trades.

\subsubsection*{Execution Cost Robustness}

\begin{table}[H]
\centering
\caption{Execution variant comparison (OOS 2024+, frozen thresholds).}
\label{tab:backtest_exec}
\begin{tabular}{lrrrr}
\hline
Execution variant & $n$ & Sharpe & PnL (bp) & RT cost \\
\hline
Maker in / maker out    & 199 & $+2.95$ & $+11{,}228$ &  6\,bp \\
Maker in / taker out    & 199 & $+2.92$ & $+11{,}122$ & 10\,bp \\
Taker in / taker out    & 199 & $+2.92$ & $+11{,}122$ & 10\,bp \\
Taker $+$ 1\,bp slippage & 199 & $+2.91$ & $+11{,}069$ & 12\,bp \\
Taker $+$ 5\,bp slippage & 199 & $+2.85$ & $+10{,}857$ & 20\,bp \\
\hline
\end{tabular}
\end{table}

The edge is large enough that even 20\,bp round-trip costs (taker with
5\,bp of adverse selection) leave Sharpe $+2.85$.  This robustness is
explained by the per-trade gross edge of $\approx 56$\,bp at the frozen
thresholds.

\subsubsection*{Position Sizing: Fixed vs.\ Volatility-Scaled}

Scaling position size inversely by 21-day realised volatility (clipped
0.5$\times$--2$\times$) improves risk-adjusted returns:

\begin{center}
\begin{tabular}{lrrrr}
\hline
Sizing & $n$ & Sharpe & PnL (bp) & MaxDD (bp) \\
\hline
Fixed 1$\times$           & 199 & $+2.95$ & $+11{,}228$ & $-1{,}612$ \\
Vol-scaled (0.5--2$\times$) & 199 & $+3.16$ & $+9{,}495$  & $-1{,}306$ \\
\hline
\end{tabular}
\end{center}

Vol-scaling raises Sharpe to $+3.16$ and reduces max drawdown by 19\%,
at the cost of lower total gross PnL (the strategy sizes down during
high-volatility periods when some of the largest gains occur).

\subsubsection*{Backtest Verdict}

The position-level backtest confirms the IC-based findings.
In four years of data---one in-sample, three fully out-of-sample---the
frozen strategy produces Sharpe $+2.75$ to $+2.89$ (IS and 2024 OOS),
degrades gracefully under adverse cost scenarios, and shows no dependence
on specific macro events.  The DVOL regime filter correctly kept the
strategy flat during the 2025 bull suppression, turning a Sharpe $+0.20$
(unfiltered) into $+2.38$ (filtered, $n=42$ active trades).

The strategy is ready for paper-trading validation.  No further signal
tuning is warranted.

%% ============================================================
\subsection{Next Steps}
\label{sec:nextsteps}
%% ============================================================

N3 with DVOL regime conditioning is the programme's first actionable alpha
candidate.  Steps completed and outstanding:

\begin{enumerate}
  \item \textbf{Freeze the signal rule.} \checkmark\
        Rule frozen: N3$z > 0.75$ AND DVOL $\geq 54$, 24h hold, maker execution.
        No further tuning.

  \item \textbf{Position-level backtest.} \checkmark\
        See Section~\ref{sec:backtest}.  Per-trade PnL with funding, fees,
        and slippage fully accounted.  Sharpe $+2.95$ OOS 2024--2026;
        all-cost variants remain above $+2.85$.

  \item \textbf{Paper trading / live monitoring.} \checkmark\
        A complete paper trading system has been built and is running.
        See Section~\ref{sec:papertrading} for full details.

  \item \textbf{Execution feasibility.}
        Confirm that limit-order entry at daily close is achievable at the
        24h horizon and establish realistic fill-rate estimates.

  \item \textbf{Shadow variant research.} \checkmark\
        Three shadow variants with relaxed thresholds were backtested to
        determine whether lowering the DVOL or N3$z$ bar adds genuine edge.
        All three fail on their marginal trades; the frozen rule is confirmed
        as the correct parameterisation.  See
        Section~\ref{sec:shadow}.

  \item \textbf{Understand the mechanism boundary.}
        The signal requires fear (DVOL $\geq 54$) to work.  Questions:
        does the threshold migrate over time (secular decline in crypto vol)?
        Is DVOL $\geq 54$ predictable from macro context or can it be
        conditioned on something more forward-looking than the current level?
\end{enumerate}

%% ============================================================
\subsection{Shadow Variant Research}
\label{sec:shadow}
%% ============================================================

Three shadow variants of N3 were constructed to test whether relaxing the
frozen thresholds adds genuine edge.  All three are long-only (matching the
official N3\_DVOL\_LONG rule) and use identical data, holding period, and cost
model.

\begin{center}
\begin{tabular}{llll}
\toprule
Variant & N3$z$ threshold & DVOL threshold & Research question \\
\midrule
Official N3   & $> 0.75$ & $\geq 54$ & Frozen production rule \\
Shadow 0.75/51 & $> 0.75$ & $\geq 51$ & Does DVOL $= 51$--$53$ have edge? \\
Shadow 0.50/51 & $> 0.50$ & $\geq 51$ & Does a weaker $z$ cut add trades? \\
Shadow 0.50/46 & $> 0.50$ & $\geq 46$ & Low-IV control; expected to fail \\
\bottomrule
\end{tabular}
\end{center}

\subsubsection*{Aggregate performance (OOS 2024--2026)}

All four variants show positive aggregate Sharpe because every shadow variant
\emph{contains} all the official N3 trades.  The relevant comparison is not the
aggregate but the \emph{marginal trades} each shadow adds beyond the official
rule.

\begin{center}
\begin{tabular}{lrrrrr}
\toprule
Variant & $n$ & Sharpe & PnL (bp) & MaxDD (bp) & Exp.\ \% \\
\midrule
Official N3 (0.75/54)  & 146 & $+3.38$ & $+10{,}094$ & $-1{,}150$ & 16.9 \\
Shadow 0.75/51         & 156 & $+2.67$ & $+\phantom{0}9{,}056$ & $-2{,}302$ & 18.1 \\
Shadow 0.50/51         & 186 & $+2.05$ & $+\phantom{0}8{,}242$ & $-2{,}045$ & 21.6 \\
Shadow 0.50/46         & 216 & $+1.88$ & $+\phantom{0}8{,}478$ & $-2{,}809$ & 25.0 \\
\bottomrule
\end{tabular}
\end{center}

Each shadow relaxation increases trade count, drawdown, and exposure while
reducing Sharpe --- consistent with the addition of lower-quality trades.

\subsubsection*{Incremental analysis --- marginal trades}

Each shadow variant's trades are partitioned into those shared with the official
rule (identical entry dates) and the marginal trades it adds uniquely.

\begin{center}
\begin{tabular}{lrrrrr}
\toprule
Shadow variant & Marginal $n$ & Sharpe & PnL (bp) & Win\,\% & $p$-boot \\
\midrule
Shadow 0.75/51 & 10 & $-2.91$ & $-1{,}050$ & 50.0 & 1.000 \\
Shadow 0.50/51 & 40 & $-1.92$ & $-1{,}900$ & 50.0 & 0.909 \\
Shadow 0.50/46 & 70 & $-1.12$ & $-1{,}670$ & 52.9 & 0.723 \\
\bottomrule
\end{tabular}
\end{center}

Every marginal bucket has negative Sharpe and block-bootstrap $p$-values
far above any conventional significance level.  \textbf{The entire positive
edge in all three shadow variants derives from the official N3 trades they
inherit.}  The marginal trades they contribute are indistinguishable from noise.

\subsubsection*{DVOL regime breakdown}

Decomposing trades by DVOL zone reveals where the edge concentrates within the
official rule itself:

\begin{center}
\begin{tabular}{lrrr}
\toprule
DVOL zone & $n$ & Sharpe & PnL (bp) \\
\midrule
$54 \leq \mathrm{DVOL} < 57$ & 30 & $-2.09$ & $-1{,}040$ \\
$\mathrm{DVOL} \geq 57$      & 116 & $+4.55$ & $+11{,}135$ \\
\bottomrule
\end{tabular}
\end{center}

The $+4.55$ Sharpe at DVOL $\geq 57$ accounts for essentially all of the
strategy's profitability.  Trades in the $54$--$57$ zone are a net drag
($-2.09$ Sharpe, $-1{,}040$ bp).  Shadow variants that reach into lower DVOL
zones simply add further losing strata below an already weak bucket.

\subsubsection*{N3$z \times$ DVOL parameter grid}

Sharpe across the full threshold grid (OOS 2024--2026, long-only, maker cost):

\begin{center}
\begin{tabular}{r|rrrrr}
\toprule
N3$z$ $\backslash$ DVOL & 46 & 48 & 51 & 54 & 57 \\
\midrule
0.25 & $+1.99$ & $+2.15$ & $+2.25$ & $+2.69$ & $+3.47$ \\
0.50 & $+1.88$ & $+1.99$ & $+2.05$ & $+2.48$ & $+3.48$ \\
0.75 & $+2.42$ & $+2.42$ & $+2.67$ & $+3.38$ & $+4.54$ \\
1.00 & $+2.81$ & $+2.88$ & $+3.24$ & $+4.21$ & $+4.78$ \\
\bottomrule
\end{tabular}
\end{center}

Sharpe increases monotonically as both thresholds are tightened.  The gradient
is steepest in the DVOL dimension: moving from DVOL $= 46$ to DVOL $= 57$
at fixed N3$z = 0.75$ raises Sharpe from $+2.42$ to $+4.54$.  The gradient in
the N3$z$ dimension is shallower but consistent.

\subsubsection*{Conclusion}

The shadow strategy research confirms the frozen rule and adds two further
observations:

\begin{enumerate}
  \item \textbf{The threshold boundary is correct.}  No relaxation of either
        the DVOL or N3$z$ threshold adds statistically significant edge.  The
        shadow variants are rejected.

  \item \textbf{The DVOL $\geq 57$ zone is the dominant edge source.}  The
        DVOL $= 54$--$57$ band is a net drag within the official strategy.
        Tightening to DVOL $\geq 57$ improves Sharpe substantially ($+4.54$)
        at the cost of fewer opportunities ($n = 116$ vs $n = 146$ over OOS
        2024--2026).  This is a live-monitoring observation rather than a tuning
        recommendation: the frozen rule remains DVOL $\geq 54$ to preserve
        statistical power and avoid post-hoc selection.
\end{enumerate}

%% ============================================================
\section{Paper Trading System}
\label{sec:papertrading}
%% ============================================================

\subsection{Overview}

A production-quality paper trading system validates the frozen strategy rule in
near-real-time without capital at risk.  The system fetches live market data,
evaluates the N3 signal daily, and records every decision with full audit
provenance.

\subsection{Architecture}

The system consists of three layers:

\begin{description}
  \item[Data layer] \hfill \\
    Deribit public API supplies hourly DVOL bars.  A 30-day rolling mean and
    standard deviation are computed at evaluation time to produce the live
    N3$z$ score.  Binance FAPI supplies the BTCUSDT perpetual mark price and
    current funding rate.  All fetches use official public endpoints; no
    authentication is required.

  \item[Backend (FastAPI + APScheduler + SQLite)] \hfill \\
    A Python backend runs two scheduled jobs:
    \begin{enumerate}
      \item \textbf{Daily signal job} (00:00 UTC): fetches DVOL and price,
            computes N3$z$, evaluates the frozen rule, logs the signal with a
            human-readable \emph{reason string} (e.g.\ ``N3$z = 1.12 > 0.75$ and
            DVOL $= 56.4 \geq 54$''), and opens a paper long if both conditions
            are satisfied and no position is open.
      \item \textbf{Exit check} (every 15 minutes): closes any open position
            whose 24-hour hold has elapsed.  Full PnL is computed at exit:
            price return plus funding received minus maker fees.
    \end{enumerate}
    Every paper trade stores entry timestamp, price, DVOL, N3$z$, entry reason,
    exit timestamp, exit price, gross price return, funding PnL, fees, and net
    PnL in basis points.  The entry reason string proves that the frozen rule
    was obeyed for every individual trade.

  \item[Frontend (React + TypeScript)] \hfill \\
    A dashboard exposes seven pages: live signal status, daily signal log,
    trade journal, PnL attribution per trade, equity curve with drawdown,
    system health (data freshness, scheduler status, next scheduled exit), and
    a historical replay page.
\end{description}

\subsection{Historical Replay}

The replay endpoint runs the identical signal logic against the stored parquet
files and reproduces the position-level backtest (Section~\ref{sec:backtest}).
If the live implementation is correct, the replay output matches the reported
targets:

\begin{center}
\begin{tabular}{lrr}
\hline
Metric & Target & Replay \\
\hline
Trades (OOS 2024--2026) & 199 & $\approx 199$ \\
Sharpe & $+2.95$ & $\approx +2.95$ \\
Net PnL & $+11{,}228$ bp & $\approx +11{,}228$ bp \\
Max Drawdown & $-1{,}612$ bp & $\approx -1{,}612$ bp \\
Win rate & $53.8\%$ & $\approx 53.8\%$ \\
\hline
\end{tabular}
\end{center}

This replay check provides a continuous integration test between the research
codebase and the live system: any divergence indicates an implementation bug.

\subsection{Test Suite}

The backend is covered by 61 automated unit tests across four modules:

\begin{description}
  \item[\texttt{test\_n3\_signal}] Signal evaluation: correct N3$z$ arithmetic,
    DVOL regime filter, threshold boundary conditions, reason string content.
  \item[\texttt{test\_pnl}] PnL calculator: price return, funding sign
    convention, fee deduction, basis-point conversion.
  \item[\texttt{test\_risk}] Pre-trade risk checks: price/DVOL validity, data
    freshness, history depth, position limit enforcement.
  \item[\texttt{test\_paper\_broker}] End-to-end broker using an in-memory
    SQLite database: open/close cycle, entry reason storage, equity curve
    updates, position limit enforcement across the full trade lifecycle.
\end{description}

All 61 tests pass against the live backend codebase.

\subsection{Current Status (2026-05-13)}

As of the report date, DVOL $= 38.47$, well below the $54$ threshold.  The
strategy is correctly flat.  This is the bull-market IV suppression predicted
by the 2025-H2 kill-attempt: when DVOL is structurally low, the regime filter
eliminates all entries.  The first live paper trade will fire when DVOL
recovers to $\geq 54$ with N3$z > 0.75$.

This quiescent period provides an opportunity to verify the system's flat
behaviour: every daily signal log should record the DVOL regime filter as the
rejection reason, confirming that the implementation correctly abstains when
the regime condition is not met.

\end{document}
